\subsection{环的更换}

命题2. 设 A、B 为两个环， \(\rho: \mathbf{B} \to \mathbf{A}\) 为一个环同态，且 E（相应地，F）为右（相应地，左）A-模。则存在且仅存在一个 Z-线性映射

\[
\phi : \rho_ {*} (\mathbf {E}) \otimes_ {\mathbf {B}} \rho_ {*} (\mathbf {F}) \rightarrow \mathbf {E} \otimes_ {\mathbf {A}} \mathbf {F}\tag{6}
\]

使得对所有 \(x \in \mathbf{E}\) 与 \(y \in \mathbf{F}\) ， \(\phi\) 在元素 \(x \otimes y\) （它属于 \(\rho_{*}(\mathbf{E}) \otimes_{\mathbf{B}} \rho_{*}(\mathbf{F})\) ）处的像是元素 \(x \otimes y\) （它属于 \(\mathbf{E} \otimes_{\mathbf{A}} \mathbf{F}\) ） ；这个 \(\mathbf{Z}\) -线性映射是满射的。

我们考虑映射 \((x,y)\mapsto x\otimes y\) ，即从 \(\rho_{\star}(E)\times \rho_{*}(F)\) 到 \(E\otimes_A F\) 的映射；它是 Z-双线性的，并且对所有 \(\beta \in B\) ，根据定义有 \((x\rho (\beta))\otimes y = x\otimes (\rho (\beta)y)\) ，因此第1号的条件 (1) 成立，从而得到 \(\phi\) 的存在性与唯一性（第1号，命题1）。后一断言由元素 \(x\otimes y\) 生成 Z-模 \(E\otimes_A F\) 这一事实得出。

映射 (6) 称为典范映射。

推论。设 \(\mathfrak{a}\) 是A的一个双边理想，使得 \(\mathfrak{a}\) 包含于 E 的零化子与 F 的零化子中，从而 E（相应地 F）具有典范的左（相应地右） \((A/\mathfrak{a})\)-模结构（§ 1，第 12 号）。于是典范同态 (6)

\[
\phi : \mathbf {E} \otimes_ {\mathbf {A}} \mathbf {F} \rightarrow \mathbf {E} \otimes_ {\mathbf {A} / \mathfrak {a}} \mathbf {F}
\]

对应于典范同态 \(\rho: \mathbf{A} \to \mathbf{A} / \mathfrak{a}\) 是恒等映射。

对所有 \(\bar{\alpha} \in \mathbf{A} / \mathfrak{a}\) 、所有 \(x \in \mathbf{E}\) 以及所有 \(y \in \mathbf{F}\) ，有 \(x\bar{\alpha} = x\alpha\) （相应地， \(\bar{\alpha}y = \alpha y\) ），其中 \(\alpha\) 取遍满足 \(\rho(\alpha) = \bar{\alpha}\) 的一切元素。若 \(\mathbf{C} = \mathbf{Z}^{(\mathbb{E} \times \mathbb{F})}\) ，则 \(\mathbf{C}\) 中由元素 \((x\alpha, y) - (x, \alpha y)\) 生成的子模等于由元素 \((x\bar{\alpha}, y) - (x, \bar{\alpha}y)\) 生成的子模。

在命题2的假设和记号下，设 \(\mathbf{E}'\) 是一个右B-模， \(\mathbf{F}'\) 一个左B-模，并考虑两个半线性映射 \(u: \mathbf{E}' \to \mathbf{E}\) , \(v: \mathrm{F}' \to \mathrm{F}\) （它们相对于同态 \(\rho: \mathrm{B} \to \mathrm{A}\) ）； \(u\) （相应地 \(v\) ) 可被视为一个 B-线性映射 \(\mathrm{E}' \to \rho_*(\mathrm{E})\) （相应地 \(\mathrm{F}' \to \rho_*(\mathrm{F})\) )，从而得到一个 Z-线性映射 \(w: \mathrm{E}' \otimes_{\mathrm{B}} \mathrm{F}' \to \rho_*(\mathrm{E}) \otimes_{\mathrm{B}} \rho_*(\mathrm{F})\) ，使得 \(w(x' \otimes y') = u(x') \otimes v(y')\) ，其中 \(x' \in \mathrm{E}', y' \in \mathrm{F}'\) ；将典范映射 (6) 与这个映射复合，便得到一个 Z-线性映射 \(w': \mathrm{E}' \otimes_{\mathrm{B}} \mathrm{F}' \to \mathrm{E} \otimes_{\mathrm{A}} \mathrm{F}\) ，使得 \(w'(x' \otimes y') = u(x') \otimes v(y')\) ，其中 \(x' \in \mathrm{E}', y' \in \mathrm{F}'\) ；若不致引起混淆，这个映射通常就记作 \(u \otimes v\) 。显然， \((u, v) \mapsto u \otimes v\) 是一个 Z-双线性映射

\[
\operatorname{Hom} _ {\mathbf {B}} \left(\mathrm{E} ^ {\prime}, \rho_ {*} (\mathrm{E})\right) \times \operatorname{Hom} _ {\mathbf {B}} \left(\mathrm{F} ^ {\prime}, \rho_ {*} (\mathrm{F})\right)\rightarrow \operatorname{Hom} _ {\mathbf {Z}} \left(\mathrm{E} ^ {\prime} \otimes_ {\mathbf {B}} \mathrm{F} ^ {\prime}, \mathrm{E} \otimes_ {\mathbf {A}} \mathrm{F}\right).
\]

此外，若 C 是第三个环， \(\sigma: \mathrm{C} \to \mathrm{B}\) 是一个同态， \(\mathrm{E}''\) 是一个右C-模， \(\mathrm{F}''\) 是一个左C-模， \(u': \mathrm{E}'' \to \mathrm{E}'\) 和 \(v': \mathrm{F}'' \to \mathrm{F}'\) 是相对于 \(\sigma\) 的半线性映射，则

\[
(u \circ u ^ {\prime}) \otimes (v \circ v ^ {\prime}) = (u \otimes v) \circ (u ^ {\prime} \otimes v ^ {\prime}).
\]

\subsection{张量积上的算子；张量积作为多重模}

在第1号的假设和记号下，对于每个自同态 \(u\) （相应地 \(v\) ，即 A-模 \(E\) （相应地 \(F\) ) 的自同态）， \(u \otimes l_F\) （相应地 \(l_E \otimes v\) ）是 \(Z\) -模 \(E \otimes_A F\) 的自同态；由第2号的公式 (5) 立即得出，映射 \(u \mapsto u \otimes l_F\) （相应地 \(v \mapsto l_E \otimes v\) ）是一个环同态

\[
\operatorname{End} _ {\mathbf {A}} (\mathbf {E}) \rightarrow \operatorname{End} _ {\mathbf {Z}} (\mathbf {E} \otimes_ {\mathbf {A}} \mathbf {F})
\]

（相应地 \(\operatorname{End}_{\mathbf{A}}(\mathbf{F}) \to \operatorname{End}_{\mathbf{Z}}(\mathbf{E} \otimes_{\mathbf{A}} \mathbf{F}))\) ；此外，

\[
(u \otimes \mathrm{l} _ {\mathrm{F}}) \circ (\mathrm{l} _ {\mathrm{E}} \otimes v) = (\mathrm{l} _ {\mathrm{E}} \otimes v) \circ (u \otimes \mathrm{l} _ {\mathrm{F}}) = u \otimes v\tag{7}
\]

因此（§ 1，第 14 号） E ⊗\_A F 关于环 End\_A(E) 与 End\_A(F) 具有典范的双模结构。

既然如此，假设在 E 上给定一个 \(((\mathbf{B}_i^{\prime});\mathbf{A},(\mathbf{C}_j^{\prime}))\) -多重模结构，在 F 上给定一个 \((\mathbf{A},(\mathbf{B}_h^{\prime});(\mathbf{C}_k^{\prime}))\) -多重模结构（§ 1，第 14 号）；这等同于说，给定了环同态 \(\mathrm{B}_i^\prime \to \mathrm{End}_{\mathbf{A}}(\mathbf{E})\) , \(\mathrm{C}_j^{\prime 0} \to \mathrm{End}_{\mathbf{A}}(\mathbf{E})\) ，其像两两可交换，以及环同态 \(\mathrm{B}_h^{\prime \prime} \to \mathrm{End}_{\mathbf{A}}(\mathbf{F})\) , \(\mathrm{C}_k^{\prime \prime 0} \to \mathrm{End}_{\mathbf{A}}(\mathbf{F})\) ，其像也两两可交换。若将这些同态分别与上面定义的典范同态 \(\mathrm{End}_{\mathbf{A}}(\mathbf{E}) \to \mathrm{End}_{\mathbf{Z}}(\mathbf{E} \otimes_{\mathbf{A}} \mathbf{F})\) 和 \(\mathrm{End}_{\mathbf{A}}(\mathbf{F}) \to \mathrm{End}_{\mathbf{Z}}(\mathbf{E} \otimes_{\mathbf{A}} \mathbf{F})\) 复合，则（考虑到 (7)）可以看出，环同态

\[
\mathrm{B} _ {i} ^ {\prime} \rightarrow \operatorname{End} _ {\mathbf {Z}} (\mathrm{E} \otimes_ {\mathrm{A}} \mathrm{F}), \quad \mathrm{C} _ {j} ^ {\prime 0} \rightarrow \operatorname{End} _ {\mathbf {Z}} (\mathrm{E} \otimes_ {\mathrm{A}} \mathrm{F})
\]

\[
\mathrm{B} _ {h} ^ {\prime \prime} \rightarrow \operatorname{End} _ {\mathbf {Z}} (\mathrm{E} \otimes_ {\mathrm{A}} \mathrm{F}), \quad \mathrm{C} _ {k} ^ {\prime \prime 0} \rightarrow \operatorname{End} _ {\mathbf {Z}} (\mathrm{E} \otimes_ {\mathrm{A}} \mathrm{F})
\]

其像两两可交换；换句话说，在 \(\mathbf{E} \otimes_{\mathbf{A}} \mathbf{F}\) 上定义了一个 \(((\mathbf{B}_i'), (\mathbf{B}_h''); (\mathbf{C}_j'), (\mathbf{C}_k'')\) -多重模结构；正是这个多重模也称为 \(((\mathbf{B}_i^{\prime});\mathbf{A},(\mathbf{C}_j^{\prime}))\) -多重模 E 与 \((\mathbf{A},(\mathbf{B}_h^{\prime\prime});(\mathbf{C}_k^{\prime\prime}))\) -多重模 F 的（相对于 A 的）张量积。这个多重模是类似于第1号中所考虑的泛映射问题的解；确切地说：

命题3. 设 G 为 \(((\mathbf{B}_i'), (\mathbf{B}_h''); (\mathbf{C}_j'), (\mathbf{C}_k''))-multimodule.\) 。

\begin{enumerate}
\def\labelenumi{(\alph{enumi})}
\tightlist
\item
  设 \(g\) 是多重模 \(\mathbf{E} \otimes_{\mathbf{A}} \mathbf{F}\) 到 \(\mathbf{G}\) 的一个线性映射。映射 \(f: (x, y) \mapsto g(x \otimes y)\) （从 \(\mathbf{E} \times \mathbf{F}\) 到 \(\mathbf{G}\) ）是 \(\mathbf{Z}\) -双线性的，且满足第1号的关系 (1) 以及条件
\end{enumerate}

\[
\left\{ \begin{array}{l l} f (\mu_ {i} ^ {\prime} x, y) = \mu_ {i} ^ {\prime} f (x, y), & f (x \nu_ {j} ^ {\prime}, y) = f (x, y) \nu_ {j} ^ {\prime} \\ f (x, \mu_ {h} ^ {\prime \prime} y) = \mu_ {h} ^ {\prime \prime} f (x, y), & f (x, y \nu_ {k} ^ {\prime \prime}) = f (x, y) \nu_ {k} ^ {\prime \prime} \end{array} \right.\tag{8}
\]

对所有 \(x \in \mathbf{E}, y \in \mathbf{F}\) , \(\mu_i' \in \mathbf{B}_i'\) , \(\nu_j' \in \mathbf{C}_j'\) , \(\mu_h'' \in \mathbf{B}_h''\) , \(\nu_k'' \in \mathbf{C}_{k}''\) 成立，其中 \(i, j, h, k\) 任意。

\begin{enumerate}
\def\labelenumi{(\alph{enumi})}
\setcounter{enumi}{1}
\tightlist
\item
  反之，设 \(f\) 是一个 \(\mathbf{Z}\) -双线性映射（从 \(\mathbf{E} \times \mathbf{F}\) 到 \(\mathbf{G}\) ），满足条件 (1)（第1号）和 (8)。则存在且仅存在一个线性映射 \(g\) （从多重模 \(\mathbf{E} \otimes_{\mathbf{A}} \mathbf{F}\) 到多重模 \(\mathbf{G}\) ），使得 \(f(x, y) = g(x \otimes y)\) ，其中 \(x \in \mathbf{E}, y \in \mathbf{F}\) 。
\end{enumerate}

断言 (a) 直接由 \(\mathbf{E} \otimes_{\mathbf{A}} \mathbf{F}\) 上多重模结构的定义得出，因为例如 \((x \otimes y)\nu_j' = (x\nu_j') \otimes y\) 。为证明 (b)，我们首先注意到，第1号的命题1给出了 \(\mathbf{Z}\) -线性映射 \(g\) 的存在性和唯一性，它满足 \(g(x \otimes y) = f(x, y)\) （其中 \(x \in \mathbf{E}\) , \(y \in \mathbf{F}\) ）；只需验证 \(g\) 关于多重模结构是线性的。由于元素 \(x \otimes y\) 生成 \(\mathbf{Z}\) -模 \(\mathbf{E} \otimes_{\mathbf{A}} \mathbf{F}\) ，只需验证关系 \(g(\mu'(x \otimes y)) = \mu_i' g(x \otimes y)\) 及其类似关系；但这由公式 \(g(x \otimes y) = f(x, y)\) 以及关系式 (8) 立即得出。

注释。\(\mathbf{E} \otimes_{\mathbf{A}} \mathbf{F}\) 的元素一般说来可以有多种方式写成 \(\sum_{i} (x_i \otimes y_i)\) 的形式，其中 \(x_i \in \mathbf{E}\) 且 \(y_i \in \mathbf{F}\) ；但要定义线性映射 \(g\) （从多重模 \(\mathbf{E} \otimes_{\mathbf{A}} \mathbf{F}\) 到一个多重模 \(\mathbf{G}\) ），无需验证：若 \(\sum_{i} (x_i \otimes y_i) = \sum_{j} (x_j' \otimes y_j')\) ，则 \(\sum_{i} g(x_i \otimes y_i) = \sum_{j} g(x_j' \otimes y_j')\) ；只需给定 \(g(x \otimes y)\) （对于 \(x \in \mathbf{E}\) 与 \(y \in \mathbf{F}\) ），并验证 \((x, y) \mapsto g(x \otimes y)\) 是 \(\mathbf{Z}\) -双线性的且满足条件 (1)（第1号）和 (8)。

设 \(\mathbf{E}'\) 是 \(((\mathbf{B}_i'),\mathbf{A},(\mathbf{C}_j')\) -多重模， \(\mathbf{F}'\) 一个 \((\mathbf{A},(\mathbf{B}_h'');(\mathbf{C}_k'')\) -多重模，且 \(u:\mathbf{E}\to \mathbf{E}',v:\mathbf{F}\to \mathbf{F}'\) 为多重模的线性映射；由定义（第2号）立即得出， \(u\otimes v\) 是从多重模 \(\mathbf{E}\otimes_{\mathbf{A}}\mathbf{F}\) 到多重模 \(\mathbf{E}'\otimes_{\mathbf{A}}\mathbf{F}'\) 的线性映射。

设 E 始终表示右 A-模，令 \(_{s}A_{d}\) 表示把环 A 视为 (A, A)-双模（ \(\S\) 1，第 14 号，例 1）；由上述可知，张量积 \(\mathbf{E} \otimes_{\mathbf{A}} (\mathbf{s}\mathbf{A}_{d})\) 具有典范的右 A-模结构，使得 \((x \otimes \lambda)\mu = x \otimes (\lambda\mu)\) ，其中 \(x \in E\) , \(\lambda \in A\) , \(\mu \in A\) 。映射 \((x, \lambda) \mapsto x\lambda\) （从 \(\mathbf{E} \times (\mathbf{s}\mathbf{A}_{d})\) 到 E ）是 Z-双线性的，并满足条件 (1)（第 1 号）和 (8)

（在后一种情形中， \(\mathbf{B}_i'\) , \(\mathbf{C}_j'\) 和 \(\mathbf{B}_h''\) 不出现，而族 \((\mathbf{C}_k'')\) 就化为 \(\mathbf{A}\) ）；因此（命题3），存在一个 A-线性映射 \(g\) （称为典范映射），它把 \(\mathbf{E} \otimes_{\mathbf{A}} (s\mathbf{A}_d)\) 映入 \(\mathbf{E}\) ，使得 \(g(x \otimes \lambda) = x\lambda\) ，其中 \(x \in \mathbf{E}\) , \(\lambda \in \mathbf{A}\) 。

命题4. 若 E 是一个右 A-模，则映射 \(h: x \mapsto x \otimes 1\) （它把 E 映入 \(\mathbf{E} \otimes_{\mathbf{A}} (s\mathbf{A}_d)\) ）是一个右 A-模同构，其逆同构 \(g\) 满足 \(g(x \otimes \lambda) = x\lambda\) ，其中 \(x \in \mathbf{E}, \lambda \in \mathbf{A}\) 。

若 \(g\) 是典范映射，则 \(g \circ h\) 是恒等映射 \(1_{\mathbf{E}}\) ，且 \(h \circ g\) 与 \(\mathbf{E} \otimes_{\mathbf{A}} (s\mathbf{A}_d)\) 到自身的恒等映射在形如 \(x \otimes y\) 的元素上一致，而这些元素生成这个 \(\mathbf{Z}\) -模；由此得出结论。

我们通常将写成 \(E \otimes_{A} A\) 而不是 \(E \otimes_{A} (s A_{d})\) ，并经常通过上述典范同构把 \(E \otimes_{A} A\) 与 E 等同。注意，若 E 还具有与其右 A-模结构相容的（左或右）B-模结构，则 g 和 h 对于 E 和 \(E \otimes_{A} A\) 上的 B-模结构也是同构（从而也是多重模同构）。

现在设 F 是一个左 A-模； \((sA_{d}) \otimes_{A} F\) （也记作 A \(\otimes_{A} F\) ）则具有典范的左 A-模结构，并且如同命题 4 那样，可定义一个从 A \(\otimes_{A} F\) 到 F 上的典范同构，它把 \(\lambda \otimes x\) 映为 \(\lambda x\) ，其逆同构为 \(x \mapsto 1 \otimes x\) 。

特别地，存在 (A, A)-双模 \((_{s}\mathbf{A}_{d})\otimes_{\mathbf{A}}(_{s}\mathbf{A}_{d})\) 到 \({}_{s}\mathbf{A}_{d}\) 上的一个典范同构，它把 \(\lambda \otimes \mu\) 映为 \(\lambda \mu\) 。

\subsection{交换环上两个模的张量积}

设 C 为一个交换环；对于每个 C-模 E，E 上的模结构与自身相容（§ 1，第 14 号）。若 E 和 F 是两个 C-模，则第 4 号的讨论使我们能够在张量积 E ⊗c F 上定义两个 C-模结构，分别使得 γ(x ⊗ y) = (γx) ⊗ y 以及 γ(x ⊗ y) = x ⊗ (γy)；但由于根据第 1 号定义 1，在此情形下 (γx) ⊗ y = x ⊗ (γy)，这两个结构是相同的。此后当我们把 E ⊗c F 作为 C-模来讨论时，除非另有说明，均指按此方式定义的结构。典范同构

\[
\sigma : \mathbf {F} \otimes_ {\mathbf {c}} \mathbf {E} \rightarrow \mathbf {E} \otimes_ {\mathbf {c}} \mathbf {F}
\]

（第 1 号，命题 1 的推论 2）则是一个 C-模同构。

由这一定义可知，若 \((a_{\lambda})_{\lambda\in\mathbb{L}}\) （相应地 \((b_{\mu})_{\mu\in\mathbb{M}}\) ）是 C-模 E（相应地 F）的一个生成系，则 \((a_{\lambda}\otimes b_{\mu})\) 是 C-模 \(E\otimes_{C}F\) 的一个生成系；特别地，若 E 和 F 是有限生成的 C-模，则 \(E\otimes_{C}F\) 亦然。

对于每个 C-模 G，那些 Z-双线性映射 \(f\) （从 \(\mathbf{E} \times \mathbf{F}\) 到 \(\mathbf{G}\) ），若它们满足

\[
f (\gamma x, y) = f (x, \gamma y) = \gamma f (x, y) \quad \text { for } x \in \mathbf {E}, y \in \mathbf {F}, \gamma \in \mathbf {C}\tag{9}
\]

于是称为 C-双线性映射，并构成一个 C-模，记作 \(\mathcal{L}_2(\mathrm{E},\mathrm{F};\mathrm{G})\) ；命题 3（第4号）定义了一个典范的 C-模同构（参见 §1，第 14 号，注记 1）。

\[
\mathcal {L} _ {2} (\mathrm{E}, \mathrm{F}; \mathrm{G}) \rightarrow \operatorname{Hom} _ {\mathbf {c}} (\mathrm{E} \otimes_ {\mathbf {c}} \mathrm{F}, \mathrm{G}).\tag{10}
\]

设 \(E'\) , \(F'\) 是两个 C-模，且 \(u: E \to E'\) , \(v: F \to F'\) 是两个 C-线性映射；则（第4号） \(u \otimes v\) 是 \(E \otimes_{\mathbb{C}} F\) 到 \(E' \otimes_{\mathbb{C}} F'\) 的 C-线性映射。进一步，立即可以看出 \((u, v) \mapsto u \otimes v\) 是 \(\operatorname{Hom}_{\mathbb{C}}(\mathbf{E}, \mathbf{E}') \times \operatorname{Hom}_{\mathbb{C}}(\mathbf{F}, \mathbf{F}')\) 到 \(\operatorname{Hom}_{\mathbb{C}}(\mathbf{E} \otimes_{\mathbb{C}} \mathbf{F}, \mathbf{E}' \otimes_{\mathbb{C}} \mathbf{F}')\) 的 C-双线性映射；因此典范地对应于它的一个 C-线性映射，称为典范映射：

\[
\operatorname{Hom} _ {\mathbf {c}} (\mathbf {E}, \mathbf {E} ^ {\prime}) \otimes_ {\mathbf {c}} \operatorname{Hom} _ {\mathbf {c}} (\mathbf {F}, \mathbf {F} ^ {\prime}) \to \operatorname{Hom} _ {\mathbf {c}} (\mathbf {E} \otimes_ {\mathbf {c}} \mathbf {F}, \mathbf {E} ^ {\prime} \otimes_ {\mathbf {c}} \mathbf {F} ^ {\prime})\tag{11}
\]

它使每个元素 \(u \otimes v\) （它属于张量积

\[
\operatorname{Hom} _ {\mathbf {c}} (\mathbf {E}, \mathbf {E} ^ {\prime}) \otimes_ {\mathbf {c}} \operatorname{Hom} _ {\mathbf {c}} (\mathbf {F}, \mathbf {F} ^ {\prime})
\]

）对应到线性映射 \(u \otimes v\) 。注意，典范映射（11）不一定是单射也不是满射（§4，习题2）。

注记. (1) 设 A, B 为两个交换环， \(\rho:B\to A\) 为一个环同态，且 E 与 F 为两个 A-模；则第 3 号中的典范映射 (6) 是一个 B-线性映射

\[
\rho_ {*} (\mathbf {E}) \otimes \rho_ {*} (\mathbf {F}) \rightarrow \rho_ {*} (\mathbf {E} \otimes_ {\mathbf {A}} \mathbf {F}).\tag{12}
\]

\begin{enumerate}
\def\labelenumi{(\arabic{enumi})}
\setcounter{enumi}{1}
\tightlist
\item
  本号中所述内容可推广到如下情形：设 E 为右 A-模，F 为左 A-模，C 为交换环，且 \(\rho:C\to A\) 为 C 到 A 的同态，使得 \(\rho(\mathbf{C})\) 包含在 A 的中心内（参见 III，§1，第 3 号）。于是我们可以考虑 C-模 \(\rho_{*}(\mathbf{E})\) 和 \(\rho_{*}(\mathbf{F})\) ，并且关于 \(\rho\) 的假设蕴含这些 C-模结构分别与 E 和 F 上的 A-模结构相容（§1，第 14 号）。张量积 \(E\otimes_{A}F\) 因此（根据第4号）被赋予两个 C-模结构，使得 \(\gamma(x\otimes y)=(x\rho(\gamma))\otimes y\) 和
\end{enumerate}

\[
\gamma (x \otimes y) = x \otimes (\rho (\gamma) y)
\]

分别对 \(\gamma\in\mathbf{C}\) , \(x\in\mathbf{E}\) , \(y\in\mathbf{F}\) 成立，并且定义 1（第 1 号）还表明这两个结构是相同的。若 \(\mathbf{E}'\) （相应地 \(\mathbf{F}'\) ）是一个右（相应地，左）A-模，且 \(u:\mathbf{E}\to\mathbf{E}'\) , \(v:\mathbf{F}\to\mathbf{F}'\) 是两个 A-线性映射，则 \(u\otimes v:\mathbf{E}\otimes_{\mathbf{A}}\mathbf{F}\to\mathbf{E}'\otimes_{\mathbf{A}}\mathbf{F}'\) 对于刚刚定义的 C-模结构而言是 C-线性的；映射 \((u,v)\mapsto u\otimes v\) :

\[
\operatorname{Hom} _ {\mathbf {A}} (\mathrm{E}, \mathrm{E} ^ {\prime}) \times \operatorname{Hom} _ {\mathbf {A}} (\mathrm{F}, \mathrm{F} ^ {\prime}) \rightarrow \operatorname{Hom} _ {\mathbf {C}} (\mathrm{E} \otimes_ {\mathbf {A}} \mathrm{F}, \mathrm{E} ^ {\prime} \otimes_ {\mathbf {A}} \mathrm{F} ^ {\prime})
\]

是 C-双线性的（对于 \(\mathrm{Hom}_{\mathbf{A}}(\mathbf{E},\mathbf{E}^{\prime})\) 和 \(\mathrm{Hom}_{\mathbf{A}}(\mathbf{F},\mathbf{F}^{\prime})\) 上的

在 §1，第 14 号，注记 1 中定义的 C-模结构），由此我们还导出一个 C-线性映射，称为典范映射

\[
\operatorname{Hom} _ {\mathbf {A}} (\mathbf {E}, \mathbf {E} ^ {\prime}) \otimes_ {\mathbf {C}} \operatorname{Hom} _ {\mathbf {A}} (\mathbf {F}, \mathbf {F} ^ {\prime}) \rightarrow \operatorname{Hom} _ {\mathbf {C}} (\mathbf {E} \otimes_ {\mathbf {A}} \mathbf {F}, \mathbf {E} ^ {\prime} \otimes_ {\mathbf {A}} \mathbf {F} ^ {\prime}).\tag{13}
\]

\begin{enumerate}
\def\labelenumi{(\arabic{enumi})}
\setcounter{enumi}{2}
\tightlist
\item
  设 \(\mathbf{A}\) 是一个整环， \(\mathbf{K}\) 为其分式域。若 \(\mathbf{E}\) 和 \(\mathbf{F}\) 是两个向量 \(\mathbf{K}\) -空间，则典范映射
\end{enumerate}

\[
\left(\mathbf {E} _ {[ \mathbf {A} ]}\right) \otimes_ {\mathbf {A}} \left(\mathbf {F} _ {[ \mathbf {A} ]}\right)\rightarrow \mathbf {E} \otimes_ {\mathbf {K}} \mathbf {F}
\]

（第 3 号及 §1，第 13 号）是双射的。只需（第4号）证明：若 \(f\) 是 \(\mathbf{E} \times \mathbf{F}\) 到向量 K-空间 \(\mathbf{G}\) 的一个 A-双线性映射，则 \(f\) 也是 K-双线性的。现在，对所有 \(\alpha \neq 0\) （在 \(\mathbf{A}\) 中）：

\[
\alpha f (\alpha^ {- 1} x, y) = f (x, y) = \alpha f (x, \alpha^ {- 1} y)
\]

，因此

\[
f (\alpha^ {- 1} x, y) = f (x, \alpha^ {- 1} y) = \alpha^ {- 1} f (x, y)
\]

，因为 G 是一个向量 K-空间。

\subsection{E ⊗\_A F 关于正合序列的性质}

命题5. 设 E、E'、E'' 为右 A-模，F 为左 A-模，且

\[
\mathrm{E} ^ {\prime} \stackrel {u} {\longrightarrow} \mathrm{E} \stackrel {v} {\longrightarrow} \mathrm{E} ^ {\prime \prime} \longrightarrow 0\tag{14}
\]

是线性映射的一个正合列。记 \(\bar{u} = u \otimes 1_{\mathbb{F}}, \bar{v} = v \otimes 1_{\mathbb{F}}\) ，则序列

\[
\mathrm{E} ^ {\prime} \otimes_ {\mathrm{A}} \mathrm{F} \xrightarrow {\bar {u}} \mathrm{E} \otimes_ {\mathrm{A}} \mathrm{F} \xrightarrow {\bar {v}} \mathrm{E} ^ {\prime \prime} \otimes_ {\mathrm{A}} \mathrm{F} \longrightarrow 0\tag{15}
\]

是 Z-同态组成的正合列。

根据第 2 号，公式（5）， \(\bar{v} \circ \bar{u} = (v \circ u) \otimes 1_{\mathbf{F}} = 0\) ；像 \(\mathrm{H} = \bar{u} (\mathbf{E}' \otimes \mathbf{F})\) 包含于核 \(\mathbf{L} = \operatorname{Ker}(\bar{v})\) ；通过取商，我们从 \(\bar{v}\) 得到一个 \(\mathbf{Z}\) -线性映射 \(f\) ，它把余核 \(\mathbf{M} = (\mathbf{E} \otimes \mathbf{F}) / \mathrm{H}\) （ \(\bar{u}\) 的）映入 \(\mathbf{E}'' \otimes \mathbf{F}\) ；必须证明 \(f\) 是双射的，因此只需定义一个 \(\mathbf{Z}\) -线性映射 \(g: \mathbf{E}'' \otimes \mathbf{F} \to \mathbf{M}\) ，使得 \(g \circ f\) 和 \(f \circ g\) 都是恒等映射。

设 \(x'' \in E'', y \in F\) ；根据假设，存在 \(x \in E\) ，使得 \(v(x) = x''\) 。我们证明，如果 \(x_1, x_2\) 是 \(E\) 的两个元素，满足 \(v(x_1) = v(x_2) = x''\) ，且 \(\phi: E \otimes F \to M\) 是典范映射，那么 \(\phi(x_1 \otimes y) = \phi(x_2 \otimes y)\) 。只需证明若 \(v(x) = 0\) 那么 \(\phi(x \otimes y) = 0\) ，这由以下事实得出： \(x = u(x')\) （其中 \(x' \in E'\) ），因此 \(x \otimes y = u(x') \otimes y = \bar{u}(x' \otimes y) \in H\) 。若把 \((x'', y)\) 映到 \(\phi(x \otimes y)\) （对所有 \(x \in E\) 满足 \(v(x) = x''\) 而言它取唯一的值），便定义了一个从 \(E'' \times F\) 到 \(M\) 的映射；该映射是 Z-双线性的，并且满足条件 (1)（第 1 号），因为 \(v(x\lambda) = x''\lambda\) 和 \((x\lambda) \otimes y = x \otimes (\lambda y)\) （对 \(x \in E\) ）；因此存在一个 Z-线性映射 \(g\) （从 \(E'' \otimes F\) 到 \(M\) ），使得 \(g(x'' \otimes y) = \phi(x \otimes y)\) （对 \(y \in F, x \in E\) 和 \(x'' = v(x)\) ）。该定义进一步证明了 \(f \circ g\) 在

\(E'' \otimes F\) 中形如 \(x'' \otimes y\) 的元素上与恒等映射一致，因此 \(f \circ g\) 是 \(E'' \otimes F\) 的恒等映射；另一方面，对于 \(x \in E\) 和 \(y \in F\) ，根据定义有 \(f(\phi(x \otimes y)) = v(x) \otimes y\) ，因此 \(g(f(\phi(x \otimes y))) = \phi(x \otimes y)\) ，并且由于形如 \(\phi(x \otimes y)\) 的元素生成 M， \(g \circ f\) 是 M 的恒等映射。

推论. 设 F、F'、F'' 为左 A-模，E 为右 A-模，且

\[
\mathrm{F} ^ {\prime} \stackrel {s} {\longrightarrow} \mathrm{F} \stackrel {t} {\longrightarrow} \mathrm{F} ^ {\prime \prime} \longrightarrow 0\tag{16}
\]

是线性映射的一个正合列。记 \(\bar{s} = 1_{\mathbb{E}} \otimes s\) , \(\bar{t} = 1_{\mathbb{E}} \otimes t\) ，则 \(\mathbf{Z}\) -同态的序列

\[
\mathrm{E} \otimes_ {\mathrm{A}} \mathrm{F} ^ {\prime} \stackrel {\hat {s}} {\longrightarrow} \mathrm{E} \otimes_ {\mathrm{A}} \mathrm{F} \stackrel {\hat {t}} {\longrightarrow} \mathrm{E} \otimes_ {\mathrm{A}} \mathrm{F} ^ {\prime \prime} \longrightarrow 0\tag{17}
\]

是正合的。

当 E（相应地，F）被视为左（相应地，右）A \(^{0}\) -模时，F ⊗A \(^{0}\) E 与 E ⊗A F 等同起来，并且对于 F' ⊗A \(^{0}\) E 和 F'\,' ⊗A \(^{0}\) E 也有类似的等同（第 1 号，命题 1 的推论 2）；该推论随后立即由命题5得出。

注. 注意一般来说，若 \(E'\) 是右 A-模 E 的一个子模，且 \(j:E'\to E\) 是典范单射，则典范映射

\[
j \otimes 1 _ {\mathbf {F}}: \mathbf {E} ^ {\prime} \otimes \mathbf {F} \rightarrow \mathbf {E} \otimes \mathbf {F}
\]

不一定是单射。换言之，对于正合列

\[
0 \longrightarrow \mathrm{E} ^ {\prime} \stackrel {u} {\longrightarrow} \mathrm{E} \stackrel {v} {\longrightarrow} \mathrm{E} ^ {\prime \prime} \longrightarrow 0\tag{18}
\]

一般不能得出结论说序列

\[
0 \longrightarrow \mathrm{E} ^ {\prime} \otimes \mathrm{F} \stackrel {u} {\longrightarrow} \mathrm{E} \otimes \mathrm{F} \stackrel {v} {\longrightarrow} \mathrm{E} ^ {\prime \prime} \otimes \mathrm{F} \longrightarrow 0\tag{19}
\]

是正合的。

例如取A = Z，E = Z， \(E' = 2Z\) ，F = Z/2Z。由于 \(E'\) 与 E 同构， \(E' \otimes F\) 与 E \(\otimes\) F 同构，而后者本身又与 F 同构（第 4 号，命题 4）。但对所有 \(x' = 2x \in E'\) （其中 \(x \in E\) ）及所有 \(y \in F\) , \(j(x') \otimes y = (2x) \otimes y = x \otimes (2y) = 0\) ，因为 2y = 0，且 \(E' \otimes F\) 在 E \(\otimes\) F 中的典范像化为 0。

换言之，必须注意区分（对于 E 的一个子模 \(E'\) 和某个元素 \(x \in E'\) ）：元素 \(x \otimes y\) “在 \(E' \otimes F\) 中计算''与元素 \(x \otimes y\) “在 \(E \otimes F\) 中计算''（换句话说，即元素 \(j(x) \otimes y\) ）之间的差别。

我们将在后面以平坦模的名义研究这样的模F，即对于每个正合序列(18)，序列(19)都是正合的（交换代数，I，§2）。

命题6. 给定两个正合序列(14)和(16)，同态 \(v \otimes t: \mathbf{E} \otimes_{\mathbf{A}} \mathbf{F} \to \mathbf{E}'' \otimes_{\mathbf{A}} \mathbf{F}''\) 是满射且其核等于

\[
\operatorname{Im} (u \otimes \mathrm{l} _ {\mathrm{F}}) + \operatorname{Im} (\mathrm{l} _ {\mathrm{E}} \otimes s)
\]

现在 \(v \otimes t = (v \otimes 1_{\mathbb{F}''}) \circ (1_{\mathbb{E}} \otimes t)\) （第 2 号，公式（5）），因此 \(v \otimes t\) 是满射的，因为根据命题5及其推论，它是两个满射同态的复合。另一方面， \(z \in \mathbb{E} \otimes \mathbb{F}\) 属于 \(v \otimes t\) 的核，当且仅当 \((1_{\mathbb{E}} \otimes t) (z)\) 属于 \(v \otimes 1_{\mathbb{F}''}\) 的核，即（根据 (15)）属于下列映射的像：

\[
u \otimes 1 _ {F ^ {\prime \prime}}: \mathrm{E} ^ {\prime} \otimes \mathrm{F} ^ {\prime \prime} \rightarrow \mathrm{E} \otimes \mathrm{F} ^ {\prime \prime}.
\]

但作为同态 \(t: \mathbf{F} \to \mathbf{F}''\) 是满射，因此

\[
\mathbf {l} _ {\mathbf {E} ^ {\prime}} \otimes t: \mathbf {E} ^ {\prime} \otimes \mathbf {F} \rightarrow \mathbf {E} ^ {\prime} \otimes \mathbf {F} ^ {\prime \prime}
\]

由命题5的推论可知是满射的，因此关于 \(z\) 的条件归结为存在一个 \(a \in \mathbf{E}' \otimes \mathbf{F}\) ，使得

\[
(1 _ {\mathrm{E}} \otimes t) (z) = (u \otimes t) (a).
\]

设 \(b = z - (u \otimes 1_{\mathbb{F}})(a)\) ；于是 \((1_{\mathbb{E}} \otimes t)(b) = 0\) ，并且根据 (17)， \(b\) 属于 \(1_{\mathbb{E}} \otimes s\) 的像，这证明了该命题。

换言之：

推论1. 设 \(\mathbf{E}'\) 是右 A-模 \(\mathbf{E}\) 的一个子模， \(\mathbf{F}'\) 是左 A-模 \(\mathbf{F}\) 的一个子模，而 \(\operatorname{Im}(\mathbf{E}' \otimes_{\mathbf{A}} \mathbf{F})\) 和 \(\operatorname{Im}(\mathbf{E} \otimes_{\mathbf{A}} \mathbf{F}')\) 为子 \(\mathbf{Z}\) -模（ \(\mathbf{E} \otimes_{\mathbf{A}} \mathbf{F}\) 中的），即典范映射 \(\mathbf{E}' \otimes_{\mathbf{A}} \mathbf{F} \to \mathbf{E} \otimes_{\mathbf{A}} \mathbf{F}\) ，

\[
\mathbf {E} \otimes_ {\mathbf {A}} \mathbf {F} ^ {\prime} \rightarrow \mathbf {E} \otimes_ {\mathbf {A}} \mathbf {F}.
\]

则存在一个典范的Z-模同构

\[
\pi : (\mathrm{E} / \mathrm{E} ^ {\prime}) \otimes_ {\mathrm{A}} (\mathrm{F} / \mathrm{F} ^ {\prime}) \rightarrow (\mathrm{E} \otimes_ {\mathrm{A}} \mathrm{F}) / (\operatorname{Im} (\mathrm{E} ^ {\prime} \otimes_ {\mathrm{A}} \mathrm{F}) + \operatorname{Im} (\mathrm{E} \otimes_ {\mathrm{A}} \mathrm{F} ^ {\prime}))\tag{20}
\]

使得，对于 \(\xi \in \mathbf{E} / \mathbf{E}'\) , \(\eta \in \mathbf{F}\) , \(\pi (\xi \otimes \eta)\) 是由所有元素 \(x \otimes y \in \mathbf{E} \otimes_{\mathbf{A}} \mathbf{F}\) （其中 \(x \in \xi\) 且 \(y \in \eta\) ）组成的类。

注意当 E 是一个 \(((B_{i}^{\prime}); A, (C_{j}^{\prime}))\) -多重模，F 是一个 \((A, (B_{h}^{\prime\prime}); (C_{k}^{\prime\prime}))\) -多重模，且 \(E^{\prime}\) 和 \(F^{\prime}\) 分别是 E 和 F 的子多重模时，同构 (20) 是对于两边的 \(((B_{i}^{\prime}), (B_{h}^{\prime\prime}); (C_{j}^{\prime}), (C_{k}^{\prime\prime}))\) -多重模结构的同构（第 3 号）。

推论2. 设 \(\mathfrak{a}\) 是 \(\mathbf{A}\) 的一个右理想， \(\mathbf{F}\) 是一个左 A-模，而 \(\mathfrak{a}\mathbf{F}\) 是子 \(\mathbf{Z}\) -模（ \(\mathbf{F}\) 中的），由形如 \(\lambda x\) 的元素生成，其中 \(\lambda \in \mathfrak{a}\) 且 \(x \in \mathbf{F}\) 。于是存在一个典范的 \(\mathbf{Z}\) -模同构

\[
\pi : (\mathrm{A} / \mathfrak {a}) \otimes_ {\mathrm{A}} \mathrm{F} \rightarrow \mathrm{F} / \mathfrak {a} \mathrm{F}\tag{21}
\]

使得，对所有 \(\bar{\lambda} \in \mathbf{A} / \mathfrak{a}\) 以及所有 \(x \in \mathbf{F}\) , \(\pi(\bar{\lambda} \otimes x)\) 是以 \(\mathfrak{aF}\) 为模的 \(\lambda x\) 的同余类，其中 \(\lambda \in \bar{\lambda}\) 。

特别地，当 A = Z 时，可以看出对于每个整数 n 和每个 Z-模 F， \((\mathbf{Z}/n\mathbf{Z}) \otimes_{\mathbf{Z}} F\) 被典范地等同于商 Z-模 F/nF。

推论 3. 设 A 是一个交换环，a 是 A 的一个理想，E 和 F 是两个 A-模，使得 a 包含在 F 的零化子中。则 (A/a)-模 E ⊗\_A F 与 (E/aE) ⊗\_\{A/a\} F 典范同构。

F 与 \(E \otimes_{A} F\) 被 a 零化，因此具有典范的 (A/a)-模结构（§1，第 12 号），并且如果我们记 \(E' = aE\) ，则 \(\operatorname{Im}(E' \otimes_{A} F) = 0\) ；于是存在 (20) 型的典范同构，它把 \(E \otimes_{A} F\) 映到 \((\mathrm{E}/a\mathrm{E}) \otimes_{A} F\) ，而后者本身又与 \((\mathrm{E}/a\mathrm{E}) \otimes_{A/a} F\) 相同（第 3 号，命题 2 的推论）。

推论4. 设 \(\mathfrak{a}, \mathfrak{b}\) 是交换环 \(\mathbf{C}\) 中的两个理想；则 \(\mathbf{C}\) -模 \((\mathbf{C} / \mathfrak{a}) \otimes_{\mathbb{C}} (\mathbf{C} / \mathfrak{b})\) 典范同构于 \(\mathbf{C}/(\mathfrak{a} + \mathfrak{b})\) 。

\subsection{乘积与直和的张量积}

设 \((\mathbf{E}_{\lambda})_{\lambda \in L}\) 是右 A-模的一个族， \((\mathbf{F}_{\mu})_{\mu \in M}\) 是左 A-模的一个族，并考虑乘积模 \(C = \prod_{\lambda \in L} E_{\lambda}\) , \(D = \prod_{\mu \in M} F_{\mu}\) 。映射 \(((x_{\lambda}), (y_{\mu})) \mapsto (x_{\lambda} \otimes y_{\mu})\) 把 \(C \times D\) 映入乘积 Z-模

\[
\prod_ {(\lambda , \mu) \in L \times M} (E _ {\lambda} \otimes_ {A} F _ {\mu})
\]

中是 Z-双线性的，并且显然满足条件（1）（第 1 号）。因此存在（第 1 号，命题 1）一个 Z-线性映射，称为典范

\[
f: \left(\prod_ {\lambda \in L} E _ {\lambda}\right) \otimes_ {A} \left(\prod_ {\mu \in M} F _ {\mu}\right)\rightarrow \prod_ {(\lambda , \mu) \in L \times M} (E _ {\lambda} \otimes_ {A} F _ {\mu})\tag{22}
\]

，使得 \(f((x_{\lambda}) \otimes (y_{\mu})) = (x_{\lambda} \otimes y_{\mu})\) 。

当 \(C = R^{L}\) , \(D = S^{M}\) （其中 R（相应地 S）是右（相应地左）A-模）时，典范映射（22）把每个张量积 \(u \otimes v\) （其中 u 是从 L 到 R 的映射，v 是从 M 到 S 的映射）对应到映射 \((\lambda, \mu) \mapsto u(\lambda) \otimes v(\mu)\) （从 \(L \times M\) 到 \(R \otimes_{A} S\) ）；即使在这种情况下，典范映射（22）一般也既不是单射也不是满射（习题3；参见命题7的推论3）。

当 \(E_{\lambda}\) 是 \(((B_{i}^{\prime}); A, (C_{j}^{\prime}))\) -多重模而 \(F_{\mu}(A, (B_{h}^{\prime\prime}); (C_{k}^{\prime\prime}))\) -多重模时，同态（22）也是关于两边的 \(((B_{i}^{\prime}), (B_{h}^{\prime\prime}); (C_{j}^{\prime}), (C_{k}^{\prime\prime}))\) -多重模结构的一个同态。

现在考虑子模 \(\mathbf{E} = \bigoplus_{\lambda \in L} \mathbf{E}_{\lambda}\) （相应地 \(\bigoplus_{\mu \in M} F_{\mu}\) ）（属于 C（相应地 D））；典范单射 \(\mathbf{E} \to \mathbf{C}\) , \(\mathbf{F} \to \mathbf{D}\) 典范地定义了一个 Z-线性映射 \(\mathbf{E} \otimes_{\mathbf{A}} \mathbf{F} \to \mathbf{C} \otimes_{\mathbf{A}} \mathbf{D}\) ，它与映射 (22) 复合后给出一个 Z-线性映射 \(g\) （从 \(\mathbf{E} \otimes_{\mathbf{A}} \mathbf{F}\) 到 \(\prod_{\lambda, \mu} (\mathbf{E}_{\lambda} \otimes_{\mathbf{A}} F_{\mu})\) ），使得

\[
g ((x _ {\lambda}) \otimes (y _ {\mu})) = (x _ {\lambda} \otimes y _ {\mu});
\]

此外，由于族 \((x_{\lambda})\) 和 \((y_{\mu})\) 具有有限支撑，因此 \((x_{\lambda} \otimes y_{\mu})\) 亦然，从而最终 g 是一个典范同态

\[
g: \left(\bigoplus_ {\lambda \in L} E _ {\lambda}\right) \otimes_ {A} \left(\bigoplus_ {\mu \in M} F _ {\mu}\right)\rightarrow \bigoplus_ {(\lambda , \mu) \in L \times M} (E _ {\lambda} \otimes_ {A} F _ {\mu}),\tag{23}
\]

它在与 (22) 相同的条件下是一个多重模同态。

命题7. 典范映射（23）是双射的。

为证明这一点，只需定义一个 \(\mathbf{Z}\) -线性映射 \(h\) （从直和 \(\mathbf{G} = \bigoplus_{(\lambda, \mu) \in L \times M} (\mathbf{E}_{\lambda} \otimes_{\mathbf{A}} \mathbf{F}_{\mu})\) 到 \(\mathbf{E} \otimes_{\mathbf{A}} \mathbf{F}\) ），使得 \(g \circ h\) 和 \(h \circ g\) 都是恒等映射。但是，要定义一个 \(\mathbf{Z}\) -线性映射（从 \(\mathbf{G}\) 到 \(\mathbf{E} \otimes_{\mathbf{A}} \mathbf{F}\) ），只需（§ 1，第 6 号，命题 6）定义一个 \(\mathbf{Z}\) -线性映射

\[
h _ {\lambda \mu}: \mathrm{E} _ {\lambda} \otimes_ {\mathrm{A}} \mathrm{F} _ {\mu} \rightarrow \mathrm{E} \otimes_ {\mathrm{A}} \mathrm{F}
\]

对于每个有序对 \((\lambda, \mu)\) ，并且我们取 \(h_{\lambda \mu} = i_{\lambda} \otimes j_{\mu}\) ，其中 \(i_{\lambda}: \mathbf{E}_{\lambda} \to \mathbf{E}\) 和 \(j_{\mu}: \mathbf{F}_{\mu} \to \mathbf{F}\) 是典范单射。那么显然 \(h \circ g\) 与恒等映射在形如 \(\left(\sum_{\lambda} x_{\lambda}\right) \otimes \left(\sum_{\mu} y_{\mu}\right)\) 的元素上一致，这些元素生成 \(\mathbf{Z}\) -模 \(\mathbf{E} \otimes_{\mathbf{A}} \mathbf{F}\) ；类似地 \(g \circ h\) 与恒等映射在形如 \(\sum_{\lambda, \mu} (x_{\lambda} \otimes y_{\mu})\) 的元素上一致，这些元素生成 \(\mathbf{Z}\) -模 \(\mathbf{G}\) ，因为对每个有序对 \((\lambda, \mu)\) ，乘积

\[
x _ {\lambda} \otimes y _ {\mu} \quad (x _ {\lambda} \in \mathrm{E} _ {\lambda}, y _ {\mu} \in \mathrm{F} _ {\mu})
\]

生成 Z-模 \(E_{\lambda} \otimes_{A} F_{\mu}\) 。由此命题得证。

设 \(u_{\lambda}: \mathbf{E}_{\lambda} \to \mathbf{E}_{\lambda}', v_{\mu}: \mathbf{F}_{\mu} \to \mathbf{F}_{\mu}'\) 是 A-同态；显然图表

\[
\begin{array}{c} \left(\bigoplus_ {\lambda} \mathrm{E} _ {\lambda}\right) \otimes_ {\mathrm{A}} \left(\bigoplus_ {\mu} \mathrm{F} _ {\mu}\right) \longrightarrow \bigoplus_ {\lambda , \mu} (\mathrm{E} _ {\lambda} \otimes_ {\mathrm{A}} \mathrm{F} _ {\mu}) \\ (\oplus u _ {\lambda}) \otimes (\oplus v _ {\mu}) \Bigg \downarrow \qquad \qquad \qquad \qquad \qquad \qquad \qquad \qquad \qquad \qquad \qquad \qquad \qquad \qquad \qquad \qquad \qquad \qquad \qquad \oplus (u _ {\lambda} \otimes v _ {\mu}) \\ \left(\bigoplus_ {\lambda} \mathrm{E} _ {\lambda} ^ {\prime}\right) \otimes_ {\mathrm{A}} \left(\bigoplus_ {\mu} \mathrm{F} _ {\mu} ^ {\prime}\right) \longrightarrow \bigoplus_ {\lambda , \mu} (\mathrm{E} _ {\lambda} ^ {\prime} \otimes_ {\mathrm{A}} \mathrm{F} _ {\mu} ^ {\prime}) \end{array}
\]

是交换的。

推论1. 若左 A-模 F 有一个基 \((b_{\mu})_{\mu \in \mathbb{M}}\) ，则 \(\mathbf{E} \otimes_{\mathbf{A}} \mathbf{F}\) 的每个元素都可以唯一地写成形式 \(\sum_{\mu} (x_{\mu} \otimes b_{\mu})\) ，其中 \(x_{\mu} \in \mathbf{E}\) 且族 \((x_{\mu})\) 具有有限支撑。 \(\mathbf{Z}\) -模 \(\mathbf{E} \otimes_{\mathbf{A}} \mathbf{F}\) 同构于 \(\mathbf{E}^{(\mathbb{M})}\) （视为 \(\mathbf{Z}\) -模）。

基 \((b_{\mu})\) 定义了 \(\mathbf{F}\) 到 \(\bigoplus_{\mu \in \mathbb{M}}\mathrm{Ab}_{\mu}\) 的一个同构，由此（根据命题7）存在一个同构 \(\mathbf{E} \otimes_{\mathbf{A}}\mathbf{F} \to \bigoplus_{\mu \in \mathbb{M}}(\mathbf{E} \otimes_{\mathbf{A}}\mathbf{Ab}_{\mu})\) ；由于 \(\xi \mapsto \xi b_{\mu}\) 是从 \(\mathbf{A}_s\) 到 \(\mathbf{Ab}_{\mu}\) 的同构， \(x \mapsto x \otimes b_{\mu}\) 是从 \(\mathbf{E}\) 到 \(\mathbf{E} \otimes_{\mathbf{A}} (\mathbf{Ab}_{\mu})\) 的同构（根据第 4 号命题4），由此得出推论。如果 \(\mathbf{E}\) 是一个 \(((\mathbf{B}_i^{\prime}); \mathbf{A}, (\mathbf{C}_j^{\prime}))\) -多重模，则典范同构 \(\mathbf{E} \otimes_{\mathbf{A}} \mathbf{F} \to \mathbf{E}^{(\mathbf{M})}\) 是一个 \(((\mathbf{B}_i^{\prime}); (\mathbf{C}_j^{\prime}))\) -多重模同构。

特别地，若 E 也具有基 \((a_{\lambda})_{\lambda \in L}\) ，则每个 \(z \in E \otimes_{A} F\) 可以以唯一的方式写成形式 \(\sum_{\lambda, \mu} (a_{\lambda} \xi_{\lambda \mu}) \otimes b_{\mu}\) ，其中 \(\xi_{\lambda \mu}\) 属于 A（并构成一个有限支撑族）；映射 \(z \mapsto (\xi_{\lambda \mu})_{(\lambda, \mu) \in L \times M}\) 是从 \(E \otimes_{A} F\) 到 \(A^{(L \times M)}\) 的同构（对于 Z-模结构，甚至对于 A 的中心上的模结构）。更特殊地：

推论 2. 若 E 和 F 是交换环 C 上的两个自由模，且 \((a_{\lambda})\) （相应地 \((b_{\mu})\) ）是 C-模 E（相应地，F）的一组基，则 \((a_{\lambda} \otimes b_{\mu})\) 是 C-模 E \(\otimes_{C} F\) 的一组基。

出于语言习惯的滥用，基 \((a_{\lambda} \otimes b_{\mu})\) 有时称为基 \((a_{\lambda})\) 和 \((b_{\mu})\) 的张量积。

注记 (1). 设 E 是自由右 A-模，F 是自由左 A-模， \((a_{\lambda})_{\lambda \in L}\) 是 E 的一组基， \((b_{\mu})_{\mu \in M}\) 是 F 的一组基。每个元素 \(z \in E \otimes_{A} F\) 都可以唯一地写成 \(\sum_{\lambda} a_{\lambda} \otimes y_{\lambda}\) ，其中 \(y_{\lambda} \in F\) ，并且也可以唯一地写成 \(\sum_{\mu} x_{\mu} \otimes b_{\mu}\) ，其中 \(x_{\mu} \in E\) 。如果我们写 \(y_{\lambda} = \sum_{\mu} \eta_{\lambda \mu} b_{\mu}\) , \(x_{\mu} = \sum_{\lambda} a_{\lambda} \xi_{\lambda \mu}\) ，其中 \(\xi_{\lambda \mu}\) 和 \(\eta_{\lambda \mu}\) 属于 A，则 \(\xi_{\lambda \mu} = \eta_{\lambda \mu}\) （对所有 \((\lambda, \mu)\) ），因为

\[
\sum_ {\lambda} \left(a _ {\lambda} \otimes \left(\sum_ {\mu} \eta_ {\lambda \mu} b _ {\mu}\right)\right) = \sum_ {\lambda , \mu} \left((a _ {\lambda} \eta_ {\lambda \mu}) \otimes b _ {\mu}\right) = \sum_ {\mu} \left(\left(\sum_ {\lambda} a _ {\lambda} \eta_ {\lambda \mu}\right) \otimes b _ {\mu}\right).
\]

推论3. 设 \((\mathbf{E}_{\lambda})_{\lambda \in L}\) 是一族右 A-模，F 是一个自由（相应地，有限生成自由）左 A-模。则典范映射 (22)

\[
\left(\prod_ {\lambda \in L} \mathbf {E} _ {\lambda}\right) \otimes_ {\mathbf {A}} \mathbf {F} \rightarrow \prod_ {\lambda \in L} \left(\mathbf {E} _ {\lambda} \otimes_ {\mathbf {A}} \mathbf {F}\right)
\]

是单射（相应地，双射）。

若 \((b_{\mu})\) 是 \(F\) 的一个基，则 \(\left(\prod_{\lambda \in L} E_{\lambda}\right) \otimes_{A} F\) 的每个元素都可以唯一地写成 \(z = \sum_{\mu} ((x_{\lambda}^{(\mu)}) \otimes b_{\mu})\) （推论1）；说它的典范像是零，就意味着对所有 \(\lambda \in L\) , \(\sum (x_{\lambda}^{(\mu)} \otimes b_{\mu}) = 0\) ，因此 \(x_{\lambda}^{(\mu)} = 0\) （对所有 \(\lambda \in L\) 以及所有 \(\mu\) ）（推论1），从而 \(z = 0\) 。

证明当 F 具有有限基时典范映射是双射的，这件事根据命题7立即归结为以下情形：

\(F = A_{s}\) ；但此时两边都被典范地等同于 \(\prod_{\lambda \in L} E_{\lambda}\) （第 4 号，命题4），且在这些等同之后，典范映射 (22) 变为恒等映射。

推论4. 设 A 是一个无零因子的环，E 是一个自由右 A-模，F 是一个自由左 A-模。则关系 \(x \otimes y = 0\) （在 \(\mathbf{E} \otimes_{\mathbf{A}} \mathbf{F}\) 中）蕴含 \(x = 0\) 或 \(y = 0\) 。

设 \((a_{\lambda})\) 是 E 的一组基， \((b_{\mu})\) 是 F 的一组基，并令 \(x = \sum_{\lambda} a_{\lambda} \xi_{\lambda}, y = \sum_{\mu} \eta_{\mu} b_{\mu}\) ；于是 \(x \otimes y = \sum_{\lambda, \mu} ((a_{\lambda} \xi_{\lambda} \eta_{\mu}) \otimes b_{\mu})\) ，且关系 \(x \otimes y = 0\) 蕴含 \(\xi_{\lambda} \eta_{\mu} = 0\) （对每一对有序指标 \((\lambda, \mu)\) ）（推论1）。因此，若 \(x \neq 0\) ，即 \(\xi_{\lambda} \neq 0\) （对至少一个 \(\lambda\) ），则可得 \(\eta_{\mu} = 0\) （对所有 \(\mu\) ），因此 \(y = 0\) 。

推论5. 设 E 为右 A-模，F 为左 A-模，M 为 E 的子模，N 为 F 的子模。若 M 是 E 的直因子且 N 是 F 的直因子，则典范同态 \(M \otimes_{A} N \to E \otimes_{A} F\) 是单射，且 \(M \otimes_{A} N\) 在该同态下的像是 Z-模 \(E \otimes_{A} F\) 的一个直因子。

这直接由命题7得出。

注意，如果 E 是一个 \(((B_{i}^{\prime}); A, (C_{j}^{\prime}))\) -多重模，F 为 \((A, (B_{h}^{\prime\prime}); (C_{k}^{\prime\prime}))\) -多重模，且 M 和 N 是这些多重模中的直因子，则 \(M \otimes N\) 是 \(((B_{i}^{\prime}), (B_{h}^{\prime\prime}); (C_{j}^{\prime}), (C_{k}^{\prime\prime}))\) -多重模 \(E \otimes F\) 的一个直因子。

推论6. 设 P 为投射左 A-模，E、F 为两个右 A-模。对每个单同态 \(u: \mathbf{E} \to \mathbf{F}\) ，同态

\[
u \otimes \mathrm{l} _ {\mathbf {P}}: \mathrm{E} \otimes_ {\mathbf {A}} \mathrm{P} \rightarrow \mathrm{F} \otimes_ {\mathbf {A}} \mathrm{P}
\]

是单射。

存在一个左 A-模 \(\mathbf{Q}\) ，使得 \(\mathbf{L} = \mathbf{P} \oplus \mathbf{Q}\) 是自由的（§2，第 2 号，命题4），且 \(u \otimes 1_{\mathbf{L}}\) （根据命题7）等同于

\[
(u \otimes 1 _ {\mathbf {P}}) \oplus (u \otimes 1 _ {\mathbf {Q}});
\]

因此只需在 P 为自由模的情形下证明该推论（§ 1，第 6 号，命题 7 的推论 1）。同样的论证把问题归结为 \(P = A_{s}\) 的情形，这由第 4 号命题4直接得出。

推论7. 设 C 为交换环。若 E 和 F 是两个投射 C-模，则 E ⊗\_C F 是投射 C-模。

这由推论5以及两个自由C-模的张量积是自由C-模（推论2）这一事实直接得出。

注记2. 在命题7的假设下，设 \(\mathbf{E}_{\lambda}^{\prime}\) 是 \(\mathbf{E}_{\lambda}, \mathbf{F}_{\mu}^{\prime}\) 中前者的子模（后者则是 \(\mathbf{F}_{\mu}\) 的子模），并设 \(\mathbf{E}^{\prime} = \bigoplus_{\lambda \in \mathbf{L}} \mathbf{E}_{\lambda}^{\prime}, \mathbf{F}^{\prime} = \bigoplus_{\mu \in \mathbf{M}} \mathbf{F}_{\mu}^{\prime}\) 。设 \(\operatorname{Im}(\mathbf{E}^{\prime} \otimes_{\mathbf{A}} \mathbf{F}^{\prime})\)

（相应地 \(\operatorname{Im}(\mathbf{E}_{\lambda}^{\prime}\otimes_{\mathbf{A}}\mathbf{F}_{\mu}^{\prime})\) ）表示 \(\mathbf{E}^{\prime}\otimes_{\mathbf{A}}\mathbf{F}^{\prime}\) （相应地 \(\mathbf{E}_{\lambda}^{\prime}\otimes_{\mathbf{A}}\mathbf{F}_{\mu}^{\prime}\) ）在典范映射下在 \(\mathbf{E}\otimes_{\mathbf{A}}\mathbf{F}\) （相应地 \(\mathbf{E}_{\lambda}\otimes_{\mathbf{A}}\mathbf{F}_{\mu}\) ）中的像；则同构 (23) 把子 \(\mathbf{Z}\) -模

\[
\operatorname{Im} \left(\mathrm{E} ^ {\prime} \otimes_ {\mathbf {A}} \mathrm{F} ^ {\prime}\right) \quad \text { and } \quad \bigoplus_ {(\lambda , \mu) \in \mathbf {L} \times \mathbf {M}} \operatorname{Im} \left(\mathrm{E} _ {\lambda} ^ {\prime} \otimes_ {\mathbf {A}} \mathrm{F} _ {\mu} ^ {\prime}\right);
\]

这直接由图表的交换性得出。

\[
\begin{array}{c} \left(\bigoplus_ {\lambda} \mathrm{E} _ {\lambda} ^ {\prime}\right) \otimes_ {\mathrm{A}} \left(\bigoplus_ {\mu} \mathrm{F} _ {\mu} ^ {\prime}\right) \longrightarrow \left(\bigoplus_ {\lambda} \mathrm{E} _ {\lambda}\right) \otimes_ {\mathrm{A}} \left(\bigoplus_ {\mu} \mathrm{F} _ {\mu}\right) \\ \Biggl \downarrow \\ \bigoplus_ {\lambda , \mu} (\mathrm{E} _ {\lambda} ^ {\prime} \otimes_ {\mathrm{A}} \mathrm{F} _ {\mu} ^ {\prime}) \longrightarrow \bigoplus_ {\lambda , \mu} (\mathrm{E} _ {\lambda} \otimes_ {\mathrm{A}} \mathrm{F} _ {\mu}) \end{array}
\]

其中竖直箭头是典范同构。

\subsection{张量积的结合性}

命题8. 设 A、B 为两个环，E 为右 A-模，F 为 (A, B)-双模，G 为左 B-模。则 E \(\otimes_{\mathbf{A}} \mathbf{F}\) 是右 B-模，F \(\otimes_{\mathbf{B}} \mathbf{G}\) 是左 A-模，并且存在且仅存在一个 Z-线性映射

\[
\phi : (\mathrm{E} \otimes_ {\mathrm{A}} \mathrm{F}) \otimes_ {\mathrm{B}} \mathrm{G} \rightarrow \mathrm{E} \otimes_ {\mathrm{A}} (\mathrm{F} \otimes_ {\mathrm{B}} \mathrm{G})
\]

，使得 \(\phi((x \otimes y) \otimes z) = x \otimes (y \otimes z)\) （对 \(x \in \mathbf{E}\) , \(y \in \mathbf{F}\) , \(z \in \mathbf{G}\) ）；此外，这个 \(\mathbf{Z}\) -线性映射是双射的（张量积的“结合性”）。

\(\mathbf{E} \otimes_{\mathbf{A}} \mathbf{F}\) 上的右 B-模结构和 \(\mathbf{F} \otimes_{\mathbf{B}} \mathbf{G}\) 上的左 A-模结构已在第4号中定义。 \(\phi\) 的唯一性是显然的，因为元素 \((x \otimes y) \otimes z\) 生成 Z-模 \((\mathbf{E} \otimes_{\mathbf{A}} \mathbf{F}) \otimes_{\mathbf{B}} \mathbf{G}\) 。为了证明 \(\phi\) 的存在性，我们注意到，对所有 \(z \in \mathbf{G}\) , \(h_z: y \mapsto y \otimes z\) 是从左 A-模 F 到左 A-模 \(\mathbf{F} \otimes_{\mathbf{B}} \mathbf{G}\) 的 A-线性映射。我们记 \(g_z = l_{\mathbf{E}} \otimes h_z\) ，它因此是 \(\mathbf{E} \otimes_{\mathbf{A}} \mathbf{F}\) 到 \(\mathbf{E} \otimes_{\mathbf{A}} (\mathbf{F} \otimes_{\mathbf{B}} \mathbf{G})\) 的一个 Z-线性映射，并考虑映射 \((t, z) \mapsto g_z(t)\) （从 \((\mathbf{E} \otimes_{\mathbf{A}} \mathbf{F} \times \mathbf{G}\) 到 \(\mathbf{E} \otimes_{\mathbf{A}} (\mathbf{F} \otimes_{\mathbf{B}} \mathbf{G})\) ）；由于 \(h_{z+z'} = h_z + h_{z'}\) （对 \(z \in \mathbf{G}\) , \(z' \in \mathbf{G}\) ），显然上述映射是 Z-双线性的。进一步，我们证明对所有 \(\mu \in \mathbf{B}\) , \(g_{\mu z}(t) = g_z(t\mu)\) 成立；显然只需对 \(t = x \otimes y\) （其中 \(x \in \mathbf{E}\) 和 \(y \in \mathbf{F}\) ）证明即可；现在

\[
g _ {\mu z} (x \otimes y) = x \otimes (y \otimes \mu z)
\]

\[
g _ {z} ((x \otimes y) \mu) = g _ {z} (x \otimes y \mu) = x \otimes (y \mu \otimes z).
\]

命题1（第 1 号）随即证明了 Z-线性映射

\[
\phi : (\mathrm{E} \otimes_ {\mathrm{A}} \mathrm{F}) \otimes_ {\mathrm{B}} \mathrm{G} \rightarrow \mathrm{E} \otimes_ {\mathrm{A}} (\mathrm{F} \otimes_ {\mathrm{B}} \mathrm{G})
\]

，使得 \(\phi(t \otimes z) = g_z(t)\) ，因此 \(\phi((x \otimes y) \otimes z) = x \otimes (y \otimes z)\) 的存在性。类似地，一个 Z-线性映射

\[
\psi : \mathbf {E} \otimes_ {\mathbf {A}} (\mathbf {F} \otimes_ {\mathbf {B}} \mathbf {G}) \rightarrow (\mathbf {E} \otimes_ {\mathbf {A}} \mathbf {F}) \otimes_ {\mathbf {B}} \mathbf {G}
\]

定义为使 \(\psi(x \otimes (y \otimes z)) = (x \otimes y) \otimes z\) ，并且显然 \(\psi \circ \phi\) 和 \(\phi \circ \psi\) 分别是 \((\mathbf{E} \otimes_{\mathbf{A}} \mathbf{F}) \otimes_{\mathbf{B}} \mathbf{G}\) 和 \(\mathbf{E} \otimes_{\mathbf{A}} (\mathbf{F} \otimes_{\mathbf{B}} \mathbf{G})\) 的恒等映射，因为它们在这些 \(\mathbf{Z}\) -模的生成系上退化为恒等映射。

显然，如果 E 是一个 \(((C_{i}^{\prime}); A, (D_{j}^{\prime}))\) -多重模，F 是一个 \((A, (C_{h}^{\prime\prime}); B, (D_{k}^{\prime\prime}))\) -多重模，G 是一个 \((B, (C_{l}^{\prime\prime}); (D_{m}^{\prime\prime}))\) -多重模，则命题 8 中定义的典范同构是一个 \(((C_{i}^{\prime}), (C_{h}^{\prime\prime}), (C_{l}^{\prime\prime}); (D_{j}^{\prime}), (D_{k}^{\prime\prime}), (D_{m}^{\prime\prime}))\) -多重模同构。特别地，若 C 是交换环且 E、F、G 是三个 C-模，则存在典范的 C-模同构

\[
(\mathbf {E} \otimes_ {\mathbf {c}} \mathbf {F}) \otimes_ {\mathbf {c}} \mathbf {G} \rightarrow \mathbf {E} \otimes_ {\mathbf {c}} (\mathbf {F} \otimes_ {\mathbf {c}} \mathbf {G}).
\]

我们将在下面看到，在一定条件下，张量积的定义可以推广到一族多重模，这将在命题 8 的假设下特别给出一个 Z-模 \(E \otimes_{A} F \otimes_{B} G\) ，它与 \((E \otimes_{A} F) \otimes_{B} G\) 和 \(E \otimes_{A} (F \otimes_{B} G)\) 两个 Z-模中的每一个都典范同构，并且后者将与它等同。

\subsection{多重模族的张量积}

设 \((\mathbf{G}_{\lambda})_{\lambda \in L}\) 是一族 \(\mathbf{Z}\) -模；一个映射 \(u\) （从集合 \(\mathbf{G} = \prod_{\lambda \in L} \mathbf{G}_{\lambda}\) 到一个 \(\mathbf{Z}\) -模）称为多重加性的（或 \(\mathbf{Z}\) -多重线性的），如果 \((x_{\lambda}) \mapsto u((x_{\lambda}))\) 关于每个变量 \(x_{\lambda}\) 都是加性的；确切地说，这意味着，对于所有 \(\mu \in L\) 以及每个元素 \((a_{\lambda}) \in \prod_{\lambda \neq \mu} \mathbf{G}_{\lambda}\) ，把 \(\mathbf{G}\) 与 \(\mathbf{G}_{\mu} \times \prod_{\lambda \neq \mu} \mathbf{G}_{\lambda}\) 典范地等同起来，

\[
u (x _ {\mu} + y _ {\mu}, (a _ {\lambda})) = u (x _ {\mu}, (a _ {\lambda})) + u (y _ {\mu}, (a _ {\lambda})) \quad \text { for } x _ {\mu}, y _ {\mu} \text { in } G _ {\mu}.\tag{24}
\]

这特别意味着 \(u((x_{\lambda})) = 0\) （若 \(x_{\lambda}\) 中有一个为零）。

我们还考虑如下的泛映射问题，其中 \(\Sigma\) 是 Z-模结构的种类，而 \(\alpha\) -映射是将 G 映入一个 Z-模的多重加性映射。其解仍可通过考虑 Z-模 \(C = Z^{(G)}\) （即 G 的元素以 Z 中系数作形式线性组合所构成的模）以及 C 中由形如

\[
\left(x _ {\mu} + y _ {\mu}, \left(z _ {\lambda}\right) _ {\lambda \neq \mu}\right) - \left(x _ {\mu}, \left(z _ {\lambda}\right) _ {\lambda \neq \mu}\right) - \left(y _ {\mu}, \left(z _ {\lambda}\right) _ {\lambda \neq \mu}\right)
\]

（其中 \(\mu \in \mathbf{L}\) , \(x_{\mu} \in \mathrm{G}_{\mu}\) , \(y_{\mu} \in \mathrm{G}_{\mu}\) 和 \(z_{\lambda} \in \mathrm{G}_{\lambda}\) ( \(\lambda \neq \mu\) ) 都是任意的）的元素生成的子 Z-模 D 来获得。商 \(\mathbf{Z}\) -模 C/D 称为 \(\mathbf{Z}\) -模族 \((\mathrm{G}_{\lambda})_{\lambda \in L}\) 的（在 \(\mathbf{Z}\) 上的）张量积，并记作 \(\bigotimes_{\lambda \in L} \mathrm{G}_{\lambda}\) ；对每一个元素 \((x_{\lambda})_{\lambda \in L}\) （属于 G 且为 C 的典范基的一个元素）， \(\bigotimes_{\lambda \in L} x_{\lambda}\) 表示该元素在 C/D 中的典范像。由上述定义可知，映射 \(\phi: (x_{\lambda}) \mapsto \bigotimes_{\lambda \in L} x_{\lambda}\) （从 G 到 \(\bigotimes_{\lambda \in L} G_{\lambda}\) ）是 Z-多重线性的，并且对于每一个将 G 映入一个 Z-模 H 的 Z-多重线性映射 \(f\) ，存在且仅存在一个 Z-线性映射 \(g: \bigotimes_{\lambda \in L} G_{\lambda} \to H\) ，使得 \(f = g \circ \phi\) ；有序对 \(\left( \bigotimes_{\lambda \in L} G_{\lambda}, \phi \right)\) 从而解决了所讨论的泛映射问题。

设 \((\mathbf{G}_{\lambda}^{\prime})_{\lambda \in L}\) 为 \(\mathbf{Z}\) -模的另一个族，并且对于所有 \(\lambda \in L\) ，设 \(v_{\lambda}: \mathbf{G}_{\lambda} \to \mathbf{G}_{\lambda}^{\prime}\) 是 \(\mathbf{Z}\) -线性映射（换言之，即交换群的同态）。则映射

\[
(x _ {\lambda}) \mapsto \bigotimes_ {\lambda \in L} v _ {\lambda} (x _ {\lambda})
\]

（从 \(\mathbf{G}\) 到 \(\bigotimes_{\lambda \in L} \mathbf{G}_{\lambda}'\) ）是 \(\mathbf{Z}\) -多重线性的，因此典范地定义了一个 \(\mathbf{Z}\) -线性映射（从 \(\bigotimes_{\lambda \in L} \mathbf{G}_{\lambda}\) 到 \(\bigotimes_{\lambda \in L} \mathbf{G}_{\lambda}'\) ），记作 \(\bigotimes_{\lambda \in L} v_{\lambda}\) ，且使得

\[
\left(\bigotimes_ {\lambda \in L} v _ {\lambda}\right) \left(\bigotimes_ {\lambda \in L} x _ {\lambda}\right) = \bigotimes_ {\lambda \in L} v _ {\lambda} (x _ {\lambda}).\tag{25}
\]

特别地，我们考虑，对于某个 \(\mu \in \mathbf{L}\) ，一个自同态 \(\theta\) （ \(\mathbf{G}_{\mu}\) 的）；我们用 \(\tilde{\theta}\) 表示 \(\bigotimes_{\lambda \in \mathbf{L}} \mathbf{G}_{\lambda}\) 的等于 \(\bigotimes_{\lambda \in \mathbf{L}} v_{\lambda}\) 的自同态（其中 \(v_{\mu} = \theta\) ， \(v_{\lambda} = 1_{\mathbf{G}_{\lambda}}\) （对 \(\lambda \neq \mu\) ））。

然后我们假设给定了一个集合 \(\Omega\) 、一个映射

\[
c: \omega \mapsto (\rho (\omega), \sigma (\omega))
\]

（从 \(\Omega\) 到 \(\mathbf{L} \times \mathbf{L}\) ）并且，对于所有 \(\omega \in \Omega\) ，给定一个自同态 \(p_{\omega}\) （ \(\mathrm{G}_{\rho(\omega)}\) 的）以及一个自同态 \(q_{\omega}\) （ \(\mathrm{G}_{\sigma(\omega)}\) 的）；对应于它们有两个自同态 \(\tilde{p}_{\omega}\) 和 \(\tilde{q}_{\omega}\) （属于 \(\mathrm{P} = \bigotimes_{\lambda \in \mathbf{L}} \mathrm{G}_{\lambda}\) ）。设 \(\mathbf{R}\) 为子 \(\mathbf{Z}\) -模（ \(\mathbf{P}\) 中的），由自同态 \(\tilde{p}_{\omega} - \tilde{q}_{\omega}\) （当 \(\omega\) 遍历 \(\Omega\) 时）的像的并集所生成。商 \(\mathbf{Z}\) -模 \(\mathrm{P/R}\) 被称为族 \((\mathrm{G}_{\lambda})_{\lambda \in \mathbf{L}}\) 相对于 \(c, p, q\) 的张量积，并且记作 \(\bigotimes_{(c, p, q)} \mathrm{G}_{\lambda}\) ；将典范同态 \(\mathrm{P} \to \mathrm{P/R}\) 与上面定义的映射 \(\phi: \mathrm{G} \to \bigotimes_{\lambda \in \mathbf{L}} \mathrm{G}_{\lambda}\) 复合，我们得到一个 \(\mathbf{Z}\) -多重线性映射 \(\phi_{(c, p, q)}: \mathrm{G} \to \bigotimes_{(c, p, q)} \mathrm{G}_{\lambda}\) ，并记 \(\phi_{(c, p, q)}((x_{\lambda})) = \bigotimes_{(c, p, q)} x_{\lambda}\) 或简记为 \(\bigotimes_{(c)} x_{\lambda}\) 。由 \(\bigotimes_{(c, p, q)} x_{\lambda}\) 和 \(\phi_{(c, p, q)}\) 组成的有序对解决了以下泛映射问题：设 \(\tilde{p}_{\omega}\) （相应地 \(\tilde{q}_{\omega}\) ）表示从

\[
\left(x _ {\rho (\omega)}, \left(x _ {\lambda}\right) _ {\lambda \neq \rho (\omega)}\right) \mapsto \left(p _ {\omega} \left(x _ {\rho (\omega)}\right), \left(x _ {\lambda}\right) _ {\lambda \neq \rho (\omega)}\right)
\]

\[
\left(\text { resp. } \left(x _ {\sigma (\omega)}, \left(x _ {\lambda}\right) _ {\lambda \neq \sigma (\omega)}\right) \mapsto \left(q _ {\omega} (x _ {\sigma (\omega)}), \left(x _ {\lambda}\right) _ {\lambda \neq \sigma (\omega)}\right)\right)
\]

到自身的映射。那么 \(\Sigma\) 取为 Z-模结构的种类，而 \(\alpha\) -映射取为 \(\mathbf{Z}\) -多重线性映射 \(u\) （从 \(\mathbf{G}\) 到一个 \(\mathbf{Z}\) -模），它们进一步满足条件

\[
u \circ \tilde {p} _ {\omega} = u \circ \tilde {q} _ {\omega}\tag{26}
\]

（对所有 \(\omega \in \Omega\) ）。由上述定义，证明是显然的。

这一构造特别恢复了第 1 号中描述的 \(\mathbf{E} \otimes_{\mathbf{A}} \mathbf{F}\) 的构造：在这种情况下，我们取 \(\mathbf{L} = \{1, 2\}\) , \(\mathbf{G}_1 = \mathbf{E}\) , \(\mathbf{G}_2 = \mathbf{F}\) , \(\Omega = \mathbf{A}\) ；此外，对于所有 \(\omega \in \mathbf{A}\) ，我们必须有 \(\rho(\omega) = 1\) , \(\sigma(\omega) = 2\) , \(p_\omega\) 是自同态 \(x \mapsto x\omega\) （ \(\mathbf{Z}\) -模 \(\mathbf{E}\) 的），而 \(q_\omega\) 是自同态 \(y \mapsto \omega y\) （ \(\mathbf{Z}\) -模 \(\mathbf{F}\) 的）。

设 \((\mathbf{G}_{\lambda}^{\prime})_{\lambda \in L}\) 是 Z-模的第二个族；保持映射 c 不变，假设对所有 \(\omega \in \Omega\) 给定了一个自同态 \(p_{\omega}^{\prime}\) （ \(\mathbf{G}_{\rho(\omega)}^{\prime}\) 的）以及一个自同态 \(q_{\omega}^{\prime}\) （ \(\mathbf{G}_{\sigma(\omega)}^{\prime}\) 的）。然后对所有 \(\lambda \in L\) ，令 \(v_{\lambda}: G_{\lambda} \to G_{\lambda}^{\prime}\) 是一个 Z-线性映射，使得对所有 \(\omega \in \Omega\) ,

\[
v _ {\rho (\omega)} \circ p _ {\omega} = p _ {\omega} ^ {\prime} \circ v _ {\rho (\omega)} \quad \text { and } \quad v _ {\sigma (\omega)} \circ q _ {\omega} = q _ {\omega} ^ {\prime} \circ v _ {\sigma (\omega)}\tag{27}
\]

（换言之，对于所有 \(\lambda \in \mathbf{L}\) , \(v_{\lambda}\) 是关于 \(\mathbf{G}_{\lambda}\) （相应地 \(\mathrm{G}_{\lambda}^{\prime}\) ）上由 \(p_{\xi}\) 和 \(q_{\eta}\) （相应地 \(p_{\xi}^{\prime}\) 和 \(q_{\eta}^{\prime}\) ）（其中 \(\xi\) 和 \(\eta\) 满足 \(\rho(\xi) = \lambda\) 和 \(\sigma(\eta) = \lambda\) ）定义的作用律的一个态射）。则映射

\[
u: (x _ {\lambda}) \mapsto \bigotimes_ {(c, p', q')} v _ {\lambda} (x _ {\lambda})
\]

（从 G 到 \(\bigotimes_{(c, p', q')} \mathrm{G}_{\lambda}'\) ）是 Z-多重线性的，并且满足条件 (26)，因此定义了一个 Z-线性映射（从 \(\bigotimes_{(c, p, q)} \mathrm{G}_{\lambda}\) 到 \(\bigotimes_{(c, p', q')} \mathrm{G}_{\lambda}'\) ），若无混淆之虞，我们将它简记为 \(\bigotimes_{(c)} v_{\lambda}\) 。

我们现在给出如此定义的一般张量积的一个“结合性”性质。设 \((\mathbf{L}_{i})_{1\leq i\leq n}\) 是 L 的一个有限划分；对每个指标 i，令 \(\Omega_{i}\) 表示 \(\Omega\) 中由满足 \(\rho(\omega)\in\mathbf{L}_{i}\) 和 \(\sigma(\omega)\in\mathbf{L}_{i}\) 的元素组成的子集；显然，诸 \(\Omega_{i}\) 两两不相交；我们设 \(\Omega^{\prime}=\Omega-\left(\bigcup_{i}\Omega_{i}\right)\) 。对每个索引 i， \(c^{(i)}\) 将表示映射 \(\omega\mapsto(\rho(\omega),\sigma(\omega))\) （从 \(\Omega_{i}\) 到 \(L_{i}\times L_{i}\) ）；对于 \(\omega\in\Omega_{i}\) ，我们记 \(p_{\omega}^{(i)}\) 和 \(q_{\omega}^{(i)}\) 以代替 \(p_{\omega}\) 和 \(q_{\omega}\) 。于是对每个 i 存在一个“部分”张量积

\[
\mathrm{F} _ {i} = \bigotimes_ {(c ^ {(i)}, p ^ {(i)}, q ^ {(i)})} \mathrm{G} _ {\lambda}.
\]

我们进一步作如下“可置换性”假设：

\begin{enumerate}
\def\labelenumi{(\Alph{enumi})}
\setcounter{enumi}{15}
\tightlist
\item
  若 \(\omega \in \Omega'\) ，则 \(p_{\omega}\) （相应地 \(q_{\omega}\) ）与自同态 \(p_{\xi}\) 和 \(q_{\eta}\) （ \(\mathrm{G}_{\rho(\omega)}\) （相应地 \(\mathrm{G}_{\sigma(\omega)}\) ）的，其中 \(\xi \notin \Omega'\) , \(\eta \notin \Omega'\) 且 \(\rho(\omega) = \rho(\xi) = \sigma(\eta)\) （相应地 \(\sigma(\omega) = \rho(\xi) = \sigma(\eta)\) ））都可交换。
\end{enumerate}

对于每个 \(\omega \in \Omega'\) ，设 \(i\) 是使得 \(\rho(\omega) \in \mathbf{L}_i\) 的索引；然后考虑族 \((v_{\lambda})_{\lambda \in \mathbf{L}_i}\) （其中 \(v_{\rho(\omega)} = p_{\omega}\) ， \(v_{\lambda} = 1_{\mathbf{G}_{\lambda}}\) （对 \(\lambda \neq \rho(\omega)\) ））；假设 (P) 意味着族 \((v_{\lambda})\) 满足条件 (27)（其中 \(p'\) 和 \(p\) 必须换成 \(p^{(i)}\) , \(q'\) 和 \(q\) 换成 \(q^{(i)}\) , \(\omega\) 换成某个元素 \(\xi\) （它遍历 \(\Omega_i\) ））；由此导出一个自同态 \(\bigotimes_{(c^{(i)})} v_{\lambda} = r_{\omega}\) （ \(\mathbf{Z}\) -模 \(\mathbf{F}_i\) 的）。类似地，一个自同态 \(s_{\omega}\) （ \(\mathbf{Z}\) -模 \(\mathbf{F}_j\) 的）由 \(q_{\omega}, j\) 定义，后者是使得 \(\sigma(\omega) \in \mathbf{L}_j\) 的索引；最后令 \(d(\omega) = (i,j)\) 。于是我们可以定义张量积 \(\bigotimes_{(d,r,s)} \mathbf{F}_i\) 以及相应的典范映射

\[
\phi_ {(d, r, s)}: \prod_ {i = 1} ^ {n} \mathrm{F} _ {i} \rightarrow \bigotimes_ {(d, r, s)} \mathrm{F} _ {i}.
\]

另一方面，对每个 i，典范映射

\[
\psi_ {i} = \phi_ {(c ^ {(i)}, p ^ {(i)}, q ^ {(i)})}: \prod_ {\lambda \in L _ {i}} G _ {\lambda} \rightarrow F _ {i};
\]

利用集合乘积的结合性，可导出一个 Z-多重线性映射 \(\psi = \phi_{(d,r,s)}\circ (\psi_i)\) （从 G 到 \(\bigotimes_{(d,r,s)}\mathrm{F}_i\) ）。我们证明有序对 \(\left(\bigotimes_{(d,r,s)}\mathrm{F}_i,\psi\right)\) 是与 \(\left(\bigotimes_{(c,p,q)}\mathrm{G}_{\lambda},\phi_{(c,p,q)}\right)\) 相同的泛映射问题的解，由此将推出存在唯一的 Z-模同构

\[
\theta \colon \bigotimes_ {(c, p, q)} \mathrm{G} _ {\lambda} \rightarrow \bigotimes_ {(d, r, s)} \mathrm{F} _ {i}
\]

，使得 \(\psi = \theta \circ \phi_{(c,p,q)}\) （集合论，IV，§ 3，第 1 号）。

对 \(n\) 作归纳，证明归结为情形 \(n = 2\) ；为简单起见，我们用 \(\mathbf{F}_1 \otimes_{(d)} \mathbf{F}_2\) 和 \(y_1 \otimes_{(d)} y_2\) 分别代替 \(\bigotimes_{(d, r, s)} \mathbf{F}_i\) 和 \(\bigotimes_{(d, r, s)} y_i\) 。考虑从 \(G\) 到 \(\mathbf{F}_1 \otimes_{(d)} \mathbf{F}_2\)

\[
h \colon (x _ {\lambda}) \to \left(\bigotimes_ {(c ^ {(1)})} x _ {\lambda}\right) \underset {(d)} {\otimes} \left(\bigotimes_ {(c ^ {(2)})} x _ {\lambda}\right).
\]

的映射。它显然是 Z-多重线性的；我们证明它对所有 \(\omega\in\Omega\) 满足条件（26）。当 \(\omega\in\Omega_{1}\) 或 \(\omega\in\Omega_{2}\) 时这是显然的；否则，为固定想法起见，假设 \(\rho(\omega)\in\mathbf{L}_{1}\) 和 \(\sigma(\omega)\in\mathbf{L}_{2}\) ，则 \(h\circ\tilde{p}_{\omega}\) 和 \(h\circ\tilde{q}_{\omega}\) 在 \((x_{\lambda})\) 分别为

\[
\left(r _ {\omega} \left(\bigotimes_ {(c ^ {(1)})} x _ {\lambda}\right)\right) \underset {(d)} {\otimes} \left(\bigotimes_ {(c ^ {(2)})} x _ {\lambda}\right) \quad \text { and } \quad \left(\bigotimes_ {(c ^ {(1)})} x _ {\lambda}\right) \underset {(d)} {\otimes} \left(s _ {\omega} \left(\bigotimes_ {(c ^ {(2)})} x _ {\lambda}\right)\right)
\]

处的值，根据 \(\mathbf{F}_1\otimes_{(d)}\mathbf{F}_2\)

的定义也相等。既然如此，设 \(u\) 是一个 \(\mathbf{Z}\) -多重线性映射（从 \(G\) 到一个 \(\mathbf{Z}\) -模 \(H\) ），满足条件（26）；我们将定义一个 \(\mathbf{Z}\) -线性映射 \(v: F_1 \otimes F_2 \to H\) ，使得 \(u = v \circ h\) ，而这将证明我们的断言（重复第 1 号中命题 1 的推论 1 的论证）。对所有 \(z_2 = (x_\lambda)_{\lambda \in L_2}\) ，我们考虑从 \(\prod_{\lambda \in L_1} G_\lambda\) 到 \(H\) 的“部分”线性映射

\[
u \left(\cdot , z _ {2}\right): \left(x _ {\lambda}\right) _ {\lambda \in L _ {1}} \mapsto u \left(\left(x _ {\lambda}\right) _ {\lambda \in L _ {1}}, z _ {2}\right) = u \left(\left(x _ {\lambda}\right) _ {\lambda \in L}\right).\tag{28}
\]

显然它是 Z-多重线性的，并且对 \(\omega\in\Omega_{1}\) 满足条件（26）；因此根据定义，存在一个 Z-线性映射 \(y_{1}\mapsto w_{1}(y_{1},x_{2})\) （从 \(F_{1}\) 到 H 中），使得

\[
w _ {1} \left(\bigotimes_ {(c ^ {(1)})} x _ {\lambda}, z _ {2}\right) = u \left(\left(x _ {\lambda}\right) _ {\lambda \in L _ {1}}, z _ {2}\right)\tag{29}
\]

接下来我们考虑映射

\[
u _ {2}: (x _ {\lambda}) _ {\lambda \in L _ {2}} \mapsto w _ {1} (\cdot , (x _ {\lambda}) _ {\lambda \in L _ {2}})
\]

（一个从 \(\prod_{\lambda \in L_2} G_\lambda\) 到 \(\mathrm{Hom}_{\mathbf{Z}}(\mathrm{F}_1, \mathrm{H})\) 中的映射）；它显然是 \(\mathbf{Z}\) -多重线性的，并且对 \(\omega \in \Omega_2\) 满足条件 (26)，这可由关于 \(u\) 的假设以及关系式 (28) 和 (29)、并注意到形如 \(\bigotimes_{(c^{(1)})} x_\lambda\) 的元素生成 \(\mathbf{Z}\) -模 \(\mathrm{F}_1\) 这一事实得出。因此存在一个 \(\mathbf{Z}\) -线性映射

\[
w _ {2}: \mathrm{F} _ {2} \rightarrow \operatorname{Hom} _ {\mathbf {Z}} (\mathrm{F} _ {1}, \mathrm{H})
\]

，使得

\[
w _ {2} \left(\bigotimes_ {(c ^ {(2)})} x _ {\lambda}\right) = u _ {2} \left(\left(x _ {\lambda}\right) _ {\lambda \in L _ {2}}\right)
\]

或等价地

\[
\left(w _ {2} \left(\bigotimes_ {(c ^ {(2)})} x _ {\lambda}\right)\right) \left(\bigotimes_ {(c ^ {(1)})} x _ {\lambda}\right) = u \left(\left(x _ {\lambda}\right) _ {\lambda \in L}\right).\tag{30}
\]

我们现在考虑，对于 \(y_{1} \in F_{1}, y_{2} \in F_{2}\) ，H的元素

\[
w (y _ {1}, y _ {2}) = (w _ {2} (y _ {2})) (y _ {1}).\tag{31}
\]

显然 \(w\) 是一个 \(\mathbf{Z}\) -双线性映射（从 \(\mathbf{F}_1 \times \mathbf{F}_2\) 到 \(H\) ）。我们进一步证明，对所有 \(\omega \in \Omega'\) ，（为固定思路，设 \(\rho(\omega) \in L_1\) 和 \(\sigma(\omega) \in L_2\) ）

\[
w (r _ {\omega} (y _ {1}), y _ {2}) = w (y _ {1}, s _ {\omega} (y _ {2})).\tag{32}
\]

只需验证此关系当 \(y_{1}\) （相应地 \(y_{2}\) ）具有形式 \(\bigotimes_{(c^{(1)})} x_{\lambda}\) （相应地 \(\bigotimes_{(c^{(2)})} x_{\lambda}\) ）时成立，因为这些元素生成 Z-模 \(F_{1}\) （相应地 \(F_{2}\) ）。但根据定义， \(r_{\omega}\left(\bigotimes_{(c^{(1)})} x_{\lambda}\right) = \bigotimes_{(c^{(1)})} x_{\lambda}'\) ，其中 \(x_{\rho(\omega)}' = p_{\omega}(x_{\rho(\omega)})\) ，而 \(x_{\lambda}' = x_{\lambda}\) （对于 \(\lambda \neq \rho (\omega)\) ， \(\mathbf{L}_1\) 中）；类似地 \(s_{\omega}\left(\bigotimes_{(c^{(2)})}x_{\lambda}\right) = \bigotimes_{(c^{(2)})}x_{\lambda}^{\prime \prime}\) ，其中 \(x_{\sigma (\omega)}^{\prime \prime} = q_{\omega}(x_{\sigma (\omega)})\) ，而 \(x_{\lambda}^{\prime \prime} = x_{\lambda}\) （对于 \(\lambda \neq \sigma (\omega)\) ， \(\mathbf{L}_2\) 中）；利用 (30) 和 (31)，关系式 (32) 随后由 (26) 得出。因此存在一个 Z-线性映射 \(v\) （从 \(\mathrm{F}_1\otimes \mathrm{F}_2\) 到 H 中），使得 \(v(y_{1}\otimes y_{2}) = w(y_{1},y_{2})\) ，并且由 (30) 和 (31) 可知， \(v\circ h = u\) 。

上述定义的一般张量积最重要的特殊情形如下：我们从一族环 \((\mathbf{A}_i)_{1\leqslant i\leqslant n - 1}\) 和一族 \((\mathbf{E}_i)_{1\leqslant i\leqslant n}\) 出发，其中 \(\mathbf{E}_1\) 是一个右 \(\mathbf{A}_1\) -模， \(\mathbf{E}_n\) 是一个左 \(\mathbf{A}_{n - 1}\) -模，并且对于 \(2\leqslant i\leqslant n - 1\) ， \(\mathbf{E}_i\) 是一个 \((\mathbf{A}_{i - 1},\mathbf{A}_i)\) -双模。则上述定义按如下方式应用：L 是集合 \([1,n]\) ， \(\mathrm{G}_i = \mathrm{E}_i\) ， \(\Omega\) 是诸 \(\mathbf{A}_i\) ( \(1\leqslant i\leqslant n - 1\) ) 之和所成的集合。对于 \(\omega \in \mathbf{A}_i\) ( \(1\leqslant i\leqslant n - 1\) )，取 \(\rho (\omega) = i\) ， \(\sigma (\omega) = i + 1\) ， \(p_{\omega}\) 是自同态 \(x\mapsto x\omega\) （ Z-模 \(\mathbf{E}_i\) 的）， \(q_{\omega}\) 是自同态 \(y\mapsto \omega y\) （ Z-模 \(\mathbf{E}_{i + 1}\) 的）；相应的张量积记为

\[
\mathbf {E} _ {1} \otimes_ {\mathbf {A} _ {1}} \mathbf {E} _ {2} \otimes_ {\mathbf {A} _ {2}} \mathbf {E} _ {3} \otimes \dots \otimes_ {\mathbf {A} _ {n - 2}} \mathbf {E} _ {n - 1} \otimes_ {\mathbf {A} _ {n - 1}} \mathbf {E} _ {n}\tag{33}
\]

（一种记号，其中 \(A_{i}\) 有时可省略），该张量积的元素 \(\bigotimes_{(c,p,q)} x_{i}\) ，对于一族 \((x_{i})\) （满足 \(x_{i} \in E_{i}\) ， \(1 \leqslant i \leqslant n\) ），若不会引起混淆，可以写成 \(x_{1} \otimes x_{2} \otimes \cdots \otimes x_{n}\) ；类似的记号也用于 Z-线性映射 \(\bigotimes_{(c)} v_{i}\) 。假设 (P) 对 \([1,n]\) 的每个划分都成立，因为诸 \(E_{i}\) 是双模（ \(2 \leqslant i \leqslant n - 1\) ）。当 n = 3 时，我们由此定义了第 8 号中提到的 Z-模 \(E \otimes_{A} F \otimes_{B} G\) ，并重新得到了第 8 号中的命题 8（第 8 号）。

当每个 \(E_{i}\) 都是多重模（对于 \(2 \leqslant i \leqslant n - 1\) ， \(A_{i-1}\) 是在左边作用的一个环， \(A_{i}\) 是在右边作用的一个环，且对 i = 1 和 i = n 有类似条件），如第 4 号所述，则在 \(E_{1} \otimes_{A_1} E_{2} \otimes \cdots \otimes_{A_{n-1}} E_{n}\) 上定义了关于除 \(A_{i}\) （它们在 \(E_{i}\) ( \(1 \leqslant i \leqslant n\) ) 上作用）之外的所有环的多重模结构。

特别地，设 C 为一个交换环， \((\mathbf{E}_{i})_{1\leqslant i\leqslant n}\) 为一族 C-模。通过给 \(E_{1}\) 和 \(E_{n}\) 两个与给定结构相同的 C-模结构，并给 \(E_{i}\) （对于 \(2\leqslant i\leqslant n-1\) ）三个与给定结构相同的 C-模结构，我们在张量积上定义

\[
\mathbf {E} _ {1} \otimes_ {\mathbf {c}} \mathbf {E} _ {2} \otimes_ {\mathbf {c}} \mathbf {E} _ {3} \otimes \dots \otimes_ {\mathbf {c}} \mathbf {E} _ {n - 1} \otimes_ {\mathbf {c}} \mathbf {E} _ {n}\tag{34}
\]

这样便得到 \(n\) 个彼此相容且实际上相同的 C-模结构，因为对所有 \(\gamma \in \mathbf{C}\) 和 \((x_{i})\in \prod_{i = 1}^{n}\mathbf{E}_{i}\) ，根据定义

\[
\begin{array}{r l} \left(\gamma x _ {1}\right) \otimes x _ {2} \otimes \dots \otimes x _ {n} \\ & = x _ {1} \otimes \left(\gamma x _ {2}\right) \otimes \dots \otimes x _ {n} = \dots = x _ {1} \otimes x _ {2} \otimes \dots \otimes \left(\gamma x _ {n}\right). \end{array}
\]

当我们把张量积（34）称为 C-模时，除非另有说明，我们总是指具有这种结构的张量积；张量积（34）若无混淆，也记作 \(\bigotimes_{1\leqslant i\leqslant n}\mathbf{E}_i\) 。对每个 C-模 G，从 \(\prod_{i=1}^{n}\mathbf{E}_i\) 到 G 中的 Z-多重线性映射，若对每个指标 \(i\) 满足关系

\[
f (x _ {1}, \dots , x _ {i - 1}, \gamma x _ {i}, x _ {i + 1}, \dots , x _ {n}) = \gamma f (x _ {1}, \dots , x _ {n})\tag{35}
\]

对 \(\gamma \in \mathbf{C}\) 和 \((x_{i})\in \prod_{i}\mathbf{E}_{i}\) ，则称为 C-多重线性的，并构成一个 C-模，记为 \(\mathcal{L}_n(\mathbf{E}_1,\dots ,\mathbf{E}_n;\mathbf{G})\) ；张量积（34）的泛性质随后使我们能够定义一个典范的 C-模同构

\[
\mathcal {L} _ {n} \left(\mathrm{E} _ {1}, \dots , \mathrm{E} _ {n}; \mathrm{G}\right)\rightarrow \operatorname{Hom} _ {\mathrm{C}} \left(\mathrm{E} _ {1} \otimes_ {\mathrm{C}} \mathrm{E} _ {2} \otimes \dots \otimes_ {\mathrm{C}} \mathrm{E} _ {n}, \mathrm{G}\right)\tag{36}
\]

它将每个C-多重线性映射f与满足以下条件的C-线性映射g相关联：

\[
f (x _ {1}, \dots , x _ {n}) = g (x _ {1} \otimes x _ {2} \otimes \dots \otimes x _ {n}).
\]

从 \(E_{1} \times \cdots \times E_{n}\) 到 C 中的一个 C-多重线性映射也称为 n-线性形式。

设 \((\mathbf{F}_i)_{1\leqslant i\leqslant n}\) 是另一族 C-模；对于每一个由 \(n\) 个 C-线性映射 \(u_i: \mathbf{E}_i \to \mathbf{F}_i\) 组成的系统， \(u_1 \otimes u_2 \otimes \cdots \otimes u_n\) （也记作 \(\bigotimes_{1\leqslant i\leqslant n} u_i\) ）是一个 C-线性映射，

\[
\mathbf {E} _ {1} \otimes_ {\mathbb {C}} \mathbf {E} _ {2} \otimes \dots \otimes_ {\mathbb {C}} \mathbf {E} _ {n} \quad \text { into } \mathbf {F} _ {1} \otimes_ {\mathbb {C}} \mathbf {F} _ {2} \otimes \dots \otimes_ {\mathbb {C}} \mathbf {F} _ {n}.
\]

进一步， \((u_{1},\ldots,u_{n})\mapsto u_{1}\otimes u_{2}\otimes\cdots\otimes u_{n}\) 是从 \(\prod_{i}\operatorname{Hom}_{\mathbb{C}}(\mathbf{E}_{i},\mathbf{F}_{i})\) 到

\[
\operatorname{Hom} _ {\mathbf {c}} (\mathbf {E} _ {1} \otimes_ {\mathbf {c}} \mathbf {E} _ {2} \otimes \dots \otimes_ {\mathbf {c}} \mathbf {E} _ {n}, \mathbf {F} _ {1} \otimes_ {\mathbf {c}} \mathbf {F} _ {2} \otimes \dots \otimes_ {\mathbf {c}} \mathbf {F} _ {n}).
\]

中的一个 C-多重线性映射。因此，典范地对应于后一个映射的是一个 C-线性映射，称为典范映射

\[
\begin{array}{r l} & \operatorname{Hom} _ {\mathbf {c}} (\mathrm{E} _ {1}, \mathrm{F} _ {1}) \otimes_ {\mathbf {c}} \operatorname{Hom} _ {\mathbf {c}} (\mathrm{E} _ {2}, \mathrm{F} _ {2}) \otimes \dots \otimes_ {\mathbf {c}} \operatorname{Hom} _ {\mathbf {c}} (\mathrm{E} _ {n}, \mathrm{F} _ {n}) \\ & \qquad \to \operatorname{Hom} _ {\mathbf {c}} (\mathrm{E} _ {1} \otimes_ {\mathbf {c}} \mathrm{E} _ {2} \otimes \dots \otimes_ {\mathbf {c}} \mathrm{E} _ {n}, \mathrm{F} _ {1} \otimes_ {\mathbf {c}} \mathrm{F} _ {2} \otimes \dots \otimes_ {\mathbf {c}} \mathrm{F} _ {n}) \end{array}\tag{37}
\]

它推广了第 5 号中对 \(n = 2\) 的情形所定义的映射。

前面看到的一般结合律性质可以在此处特化如下。给定一个划分 \((\mathbf{J}_{k})_{1\leqslant k\leqslant m}\) （ N 的区间 \([1,n]\) 的），对每个 k，令 \(F_{k}\) 为张量积 \(E_{i_{1}}\otimes_{C}E_{i_{2}}\otimes\cdots\otimes_{C}E_{i_{r}}\) ，其中 \((i_{1},\ldots,i_{r})\) 是 \(J_{k}\) 的元素组成的严格递增序列。那么存在一个典范的 C-模同构（称为“结合同构”）

\[
\mathrm{F} _ {1} \otimes_ {\mathbf {c}} \mathrm{F} _ {2} \otimes \dots \otimes_ {\mathbf {c}} \mathrm{F} _ {m} \rightarrow \mathrm{E} _ {1} \otimes_ {\mathbf {c}} \mathrm{E} _ {2} \otimes \dots \otimes_ {\mathbf {c}} \mathrm{E} _ {n}
\]

在上述记号中，它把张量积

\[
y _ {1} \otimes y _ {2} \otimes \dots \otimes y _ {m}, \quad \text { where } y _ {k} = x _ {i _ {1}} \otimes x _ {i _ {2}} \otimes \dots \otimes x _ {i _ {r}},
\]

映射到张量积 \(x_{1} \otimes x_{2} \otimes \cdots \otimes x_{n}\) （其中 \(x_{i} \in E_{i}\) 对所有 i 成立）。

特别是，如果 \(\pi\) 是 \([1, n]\) 的一个排列，写 \(\mathbf{J}_k = \{\pi(k)\}\) （对于 \(1 \leqslant k \leqslant n\) ），我们就得到一个典范的（“交换性”）同构

\[
\mathbf {E} _ {\pi (1)} \otimes_ {\mathbf {C}} \mathbf {E} _ {\pi (2)} \otimes \dots \otimes_ {\mathbf {C}} \mathbf {E} _ {\pi (n)} \rightarrow \mathbf {E} _ {1} \otimes_ {\mathbf {C}} \mathbf {E} _ {2} \otimes \dots \otimes_ {\mathbf {C}} \mathbf {E} _ {n}
\]

它把 \(x_{\pi(1)} \otimes x_{\pi(2)} \otimes \cdots \otimes x_{\pi(n)}\) 映到 \(x_{1} \otimes x_{2} \otimes \cdots \otimes x_{n}\) 。我们常常将这些典范同构之下彼此对应的各种张量积等同起来。

对于 \(1 \leqslant i \leqslant n\) ，进一步假设 \(\mathbf{E}_i\) 有一组基 \((b_{\lambda_1}^{(i)})_{\lambda_i \in \mathbf{L}_i}\) ；通过对 \(n\) 作归纳，由第 7 号命题 7 的推论 2 可知，族 \((b_{\lambda_1}^{(1)} \otimes b_{\lambda_2}^{(2)} \otimes \cdots \otimes b_{\lambda_n}^{(n)})\) （其中 \((\lambda_1, \ldots, \lambda_n)\) 遍历 \(\prod_{1 \leqslant i \leqslant n} \mathbf{L}_i\) ）是 \(\bigotimes_{1 \leqslant i \leqslant n} \mathbf{E}_i\) 的一组基，有时称为所考虑的基 \((b_{\lambda_1}^{(1)})\) 的张量积。

注记. (1) 上述关于交换环上模情形的注记，可以如第 5 号注记 2 那样推广到存在张量积 \(\mathbf{E}_1 \otimes_{\mathbf{A}_1} \mathbf{E}_2 \otimes \cdots \otimes_{\mathbf{A}_{n-1}} \mathbf{E}_n\) 的情形，其中环 \(\mathbf{A}_i\) 不一定是交换的，并且对每个 \(i\) ，存在同态 \(\rho_i: \mathbf{C} \to \mathbf{A}_i\) （同一个交换环 \(\mathbf{C}\) 的），使得：(1) \(\rho_i(\mathbf{C})\) 包含于 \(\mathbf{A}_i\) 的中心；(2) 对于 \(2 \leqslant i \leqslant n - 1\) ， \(\mathbf{E}_i\) 上利用同态 \(\rho_{i-1}\) 和 \(\rho_i\) 得到的 C-模结构是一致的。于是，在 \(\mathbf{E}_1 \otimes_{\mathbf{A}_1} \mathbf{E}_2 \otimes \cdots \otimes_{\mathbf{A}_{n-1}} \mathbf{E}_n\) 上得到一个 C-模结构，并且有类似于 (13)（第 5 号）的典范映射，这些留给读者去描述。

\begin{enumerate}
\def\labelenumi{(\arabic{enumi})}
\setcounter{enumi}{1}
\tightlist
\item
  设 A、B 为两个环，E 为右 A-模， \(E'\) 为左 A-模，F 为右 B-模， \(F'\) 为左 B-模。从 \((\mathbf{E} \otimes_{\mathbf{A}} \mathbf{E}') \times (\mathbf{F} \otimes_{\mathbf{B}} \mathbf{F}')\) 到一个 Z-模 G 中的 Z-双线性映射，与从 \(E \times E' \times F \times F'\) 到 G 中满足条件
\end{enumerate}

\[
\left\{ \begin{array}{l} f (x \lambda , x ^ {\prime}, y, y ^ {\prime}) = f (x, \lambda x ^ {\prime}, y, y ^ {\prime}) \\ f (x, x ^ {\prime}, y \mu , y ^ {\prime}) = f (x, x ^ {\prime}, y, \mu y ^ {\prime}) \end{array} \right.\tag{38}
\]

（对于 \(\lambda\in\mathbf{A}\) , \(\mu\in\mathbf{B}\) , \(x\in\mathbf{E}\) , \(x'\in\mathbf{E}'\) , \(y\in\mathbf{F}\) , \(y'\in\mathbf{F}'\) ）的 Z-多重线性映射 f 成一一对应。本号中给出的一般构造将此证明归结为定义一个典范的 \(\mathbf{Z}\) -模同构（介于 \((\mathbf{E}\otimes_{\mathbf{A}}\mathbf{E}')\otimes_{\mathbf{Z}}(\mathbf{F}\otimes_{\mathbf{B}}\mathbf{F}')\) 与 \(\mathbf{E}\otimes_{\mathbf{A}}\mathbf{E}'\otimes_{\mathbf{Z}}\mathbf{F}\otimes_{\mathbf{B}}\mathbf{F}'\) 之间），这由形如 (33) 的张量积的结合性性质得出。

\section{张量积与同态模之间的关系}

\subsection{\texorpdfstring{同构 \(\operatorname{Hom}_{\mathbf{B}}(\mathbf{E} \otimes_{\mathbf{A}} \mathbf{F}, \mathbf{G}) \to \operatorname{Hom}_{\mathbf{A}}(\mathbf{F}, \operatorname{Hom}_{\mathbf{B}}(\mathbf{E}, \mathbf{G}))\) 与 \(\operatorname{Hom}_{\mathbf{C}}(\mathbf{E} \otimes_{\mathbf{A}} \mathbf{F}, \mathbf{G}) \to \operatorname{Hom}_{\mathbf{A}}(\mathbf{E}, \operatorname{Hom}_{\mathbf{C}}(\mathbf{F}, \mathbf{G}))\)}{同构 \textbackslash operatorname\{Hom\}\_\{\textbackslash mathbf\{B\}\}(\textbackslash mathbf\{E\} \textbackslash otimes\_\{\textbackslash mathbf\{A\}\} \textbackslash mathbf\{F\}, \textbackslash mathbf\{G\}) \textbackslash 到 \textbackslash operatorname\{Hom\}\_\{\textbackslash mathbf\{A\}\}(\textbackslash mathbf\{F\}, \textbackslash operatorname\{Hom\}\_\{\textbackslash mathbf\{B\}\}(\textbackslash mathbf\{E\}, \textbackslash mathbf\{G\})) 且 \textbackslash operatorname\{Hom\}\_\{\textbackslash mathbf\{C\}\}(\textbackslash mathbf\{E\} \textbackslash otimes\_\{\textbackslash mathbf\{A\}\} \textbackslash mathbf\{F\}, \textbackslash mathbf\{G\}) \textbackslash 到 \textbackslash operatorname\{Hom\}\_\{\textbackslash mathbf\{A\}\}(\textbackslash mathbf\{E\}, \textbackslash operatorname\{Hom\}\_\{\textbackslash mathbf\{C\}\}(\textbackslash mathbf\{F\}, \textbackslash mathbf\{G\}))}}

设 E 为右 A-模，F 为左 A-模，G 为 Z-模，H 为由映射 \(f: E \times F \rightarrow G\) 组成的 Z-模，这些映射是 Z-双线性的且满足

\[
f (x \lambda , y) = f (x, \lambda y) \quad \text { for } x \in \mathrm{E}, y \in \mathrm{F}, \lambda \in \mathrm{A}.\tag{1}
\]

已经看到（§ 3，第 1 号，命题 1）存在一个典范的 Z-模同态

\[
\mathrm{H} \rightarrow \operatorname{Hom} _ {\mathbf {Z}} (\mathrm{E} \otimes_ {\mathrm{A}} \mathrm{F}, \mathrm{G}).\tag{2}
\]

另一方面，在 \(\mathrm{Hom}_{\mathbf{Z}}(\mathrm{E},\mathrm{G})\) 上已定义了左 A-模结构，在 \(\mathrm{Hom}_{\mathbf{Z}}(\mathrm{F},\mathrm{G})\) 上定义了右 A-模结构（§ 3，第 3 号）；因此我们可以考虑 Z-模 \(\mathrm{Hom}_{\mathbf{A}}(\mathrm{E},\mathrm{Hom}_{\mathbf{Z}}(\mathrm{F},\mathrm{G}))\) 和 \(\mathrm{Hom}_{\mathbf{A}}(\mathrm{F},\mathrm{Hom}_{\mathbf{Z}}(\mathrm{E},\mathrm{G}))\) 。一个映射 \(f\) （从 \(\mathbf{E} \times \mathbf{F}\) 到 \(\mathbf{G}\) 的）被典范地等同于一个从 \(\mathbf{E}\) 到集合 \(\mathbf{G}^{\mathbf{F}}\) （从 \(\mathbf{F}\) 到 \(\mathbf{G}\) 的映射所成的集合）中的映射（集合论，II，§ 5，第 2 号）；通过表达后一映射属于 \(\mathrm{Hom}_{\mathbf{A}}(\mathrm{E},\mathrm{Hom}_{\mathbf{Z}}(\mathrm{F},\mathrm{G}))\) 这一事实，我们恰好得到 \(f\) 是双加性的且满足条件 (1)；由此存在一个典范同构

\[
\mathrm{H} \rightarrow \operatorname{Hom} _ {\mathbf {A}} (\mathrm{E}, \operatorname{Hom} _ {\mathbf {Z}} (\mathrm{F}, \mathrm{G}))\tag{3}
\]

类似地，可定义一个典范同构

\[
\mathrm{H} \rightarrow \operatorname{Hom} _ {\mathbf {A}} (\mathrm{F}, \operatorname{Hom} _ {\mathbf {Z}} (\mathrm{E}, \mathrm{G})).\tag{4}
\]

现在假设 E 和 G 也具有左（相应地，右）B-模结构，并且 E 上的 A-模与 B-模结构是相容的。则 \(E \otimes_{A} F\) 典范地具有一个左（相应地，右）B-模结构（ \(\S 3\) ，第 4 号）；另一方面， \(\operatorname{Hom}_{\mathbf{B}}(\mathbf{E}, \mathbf{G})\) 具有典范的左 A-模结构（ \(\S 1\) ，第 14 号）。因此我们可以考虑Z-模 \(\operatorname{Hom}_{\mathbf{B}}(\mathbf{E} \otimes_{\mathbf{A}} \mathbf{F}, \mathbf{G})\) 和 \(\operatorname{Hom}_{\mathbf{A}}(\mathbf{F}, \operatorname{Hom}_{\mathbf{B}}(\mathbf{E}, \mathbf{G}))\) ，它们是 \(\operatorname{Hom}_{\mathbf{Z}}(\mathbf{E} \otimes_{\mathbf{A}} \mathbf{F}, \mathbf{G})\) 和 \(\operatorname{Hom}_{\mathbf{A}}(\mathbf{F}, \operatorname{Hom}_{\mathbf{Z}}(\mathbf{E}, \mathbf{G}))\) 的子模（ \(\S 2\) ，第1号，定理2）。我们考察在什么条件下，一个映射 \(f \in H\) 在同构（2）和（4）下的像分别是 \(\operatorname{Hom}_{\mathbf{B}}(\mathbf{E} \otimes_{\mathbf{A}} \mathbf{F}, \mathbf{G})\) 的一个元素和 \(\operatorname{Hom}_{\mathbf{A}}(\mathbf{E}, \operatorname{Hom}_{\mathbf{B}}(\mathbf{F}, \mathbf{G}))\) 的一个元素；在这两种情形中，我们都得到相同的条件

\[
\begin{array}{r l} f (\beta x, y) & = \beta f (x, y) \\ (\text { resp.   } f (x \beta , y) & = f (x, y) \beta) \end{array}
\]

对于 \(x \in \mathbf{E}, y \in \mathbf{F}, \beta \in \mathbf{B}\) 。

类似地，设F和G是左（相应地，右）C-模，并且F上的A-模与C-模结构是相容的。那么，一个映射 \(f \in H\) 在（2）或（3）下的像分别是 \(\operatorname{Hom}_{\mathbf{C}}(\mathbf{E} \otimes_{\mathbf{A}} \mathbf{F}, \mathbf{G})\) 或 \(\operatorname{Hom}_{\mathbf{A}}(\mathbf{E}, \operatorname{Hom}_{\mathbf{C}}(\mathbf{F}, \mathbf{G}))\) 的一个元素，当且仅当它满足相同的条件

\[
\begin{array}{c} f (x, \gamma y) = \gamma f (x, y) \\ (\text { resp. } f (x, y \gamma) = f (x, y) \gamma) \end{array}
\]

对于 \(x \in \mathbf{E}, y \in \mathbf{F}, \gamma \in \mathbf{C}\) 。

因此，我们已建立了如下结果（采用上文引入的记号）：

命题1. (a) 设E是一个(B, A)-双模，F是一个左A-模，G是一个左B-模。对于每一个映射 \(g \in \operatorname{Hom}_{\mathbf{B}}(\mathbf{E} \otimes_{\mathbf{A}} \mathbf{F}, \mathbf{G})\) ，设 \(g'\) 为从 F 到 \(\operatorname{Hom}_{\mathbf{B}}(\mathbf{E}, \mathbf{G})\) 中的映射，定义为 \((g'(y))(x) = g(x \otimes y)\) （对于 \(x \in \mathbf{E}, y \in \mathbf{F}\) ）。映射 \(g \mapsto g'\) 是一个同构

\[
\beta : \operatorname{Hom} _ {\mathbf {B}} (\mathbf {E} \otimes_ {\mathbf {A}} \mathbf {F}, \mathbf {G}) \rightarrow \operatorname{Hom} _ {\mathbf {A}} (\mathbf {F}, \operatorname{Hom} _ {\mathbf {B}} (\mathbf {E}, \mathbf {G})).\tag{5}
\]

\begin{enumerate}
\def\labelenumi{(\alph{enumi})}
\setcounter{enumi}{1}
\tightlist
\item
  设 \(\mathbf{E}\) 为一个右 \(\mathbf{A}\) -模， \(\mathbf{F}\) 一个 \((\mathbf{A}, \mathbf{C})\) -双模且 \(\mathbf{G}\) 一个右 \(\mathbf{C}\) -模。对于每个映射 \(h \in \operatorname{Hom}_{\mathbf{C}}(\mathbf{E} \otimes_{\mathbf{A}} \mathbf{F}, \mathbf{G})\) ，设 \(h'\) 是从 \(\mathbf{E}\) 到 \(\operatorname{Hom}_{\mathbf{C}}(\mathbf{F}, \mathbf{G})\) 中的映射，定义为 \((h'(x))(y) = h(x \otimes y)\) （对于 \(x \in \mathbf{E}\) , \(y \in \mathbf{F}\) ）。映射 \(h \mapsto h'\) 是一个同构
\end{enumerate}

\[
\gamma : \operatorname{Hom} _ {\mathbf {C}} (\mathrm{E} \otimes_ {\mathbf {A}} \mathrm{F}, \mathrm{G}) \rightarrow \operatorname{Hom} _ {\mathbf {A}} (\mathrm{E}, \operatorname{Hom} _ {\mathbf {C}} (\mathrm{F}, \mathrm{G})).\tag{6}
\]

特别地，B 和 C 可以取为环 A 的中心的一个子环 \(\Gamma\) ；则对于每个 \(\Gamma\) -模 G，这三个 \(\Gamma\) -模

\[
\operatorname{Hom} _ {\Gamma} (\mathrm{E} \otimes_ {\mathrm{A}} \mathrm{F}, \mathrm{G}), \quad \operatorname{Hom} _ {\mathrm{A}} (\mathrm{E}, \operatorname{Hom} _ {\Gamma} (\mathrm{F}, \mathrm{G})), \quad \operatorname{Hom} _ {\mathrm{A}} (\mathrm{F}, \operatorname{Hom} _ {\Gamma} (\mathrm{E}, \mathrm{G}))
\]

典范同构于 \(\Gamma\) -模，它由 \(\Gamma\) -双线性映射（从 \(\mathbf{E} \times \mathbf{F}\) 到 \(\mathbf{G}\) 中且满足 (1) 的）组成。更具体地说：

推论. 若 C 是交换环，且 E、F、G 是三个 C-模，则 C-模

\[
\begin{array}{l l} \operatorname{Hom} _ {\mathbf {c}} (\mathrm{E} \otimes_ {\mathbf {c}} \mathrm{F}, \mathrm{G}), & \operatorname{Hom} _ {\mathbf {c}} (\mathrm{E}, \operatorname{Hom} _ {\mathbf {c}} (\mathrm{F}, \mathrm{G})), \\ \operatorname{Hom} _ {\mathbf {c}} (\mathrm{F}, \operatorname{Hom} _ {\mathbf {c}} (\mathrm{E}, \mathrm{G})), & \mathcal {L} _ {2} (\mathrm{E}, \mathrm{F}; \mathrm{G}) \end{array}
\]

是典范同构的。
